\documentclass[final,onefignum,onetabnum]{siamart}

% ------------------------------------------------------------
% Packages (SIAM-approved)
% ------------------------------------------------------------
\usepackage{amsmath, amssymb}
\usepackage{graphicx}
\usepackage{hyperref}

% ------------------------------------------------------------
% Law Environment (SIAM-native alias)
% ------------------------------------------------------------
\newtheorem{law}{Law}
\newtheorem{axiom}{Axiom}
\newtheorem{remark}{Remark}

% ------------------------------------------------------------
% Title and Author
% ------------------------------------------------------------
\title{The Amsler Surface:\\
The Primitive Computational Power of the Superposed Hypersphere}

\author{
Aaron Schoenman\thanks{Independent Researcher, Columbus, Ohio, USA.
Email: \texttt{yummyinmytummy@naturallybaked.co.site}. ORCID: 0009-0004-8636-0927}
}

% ------------------------------------------------------------
% Begin Document
% ------------------------------------------------------------
\begin{document}
\maketitle

% ------------------------------------------------------------
% Abstract (≤250 words, one paragraph, no custom macros)
% ------------------------------------------------------------
\begin{abstract}
The classical Amsler surface is a well‑known example of a negatively curved surface governed by a nonlinear partial differential equation whose solutions exhibit characteristic oscillatory behavior. In this work, we reinterpret the Amsler surface through a geometric framework in which the surface is treated as a projection of a higher‑dimensional curvature manifold, specifically a hyperspherical structure whose collapse behavior determines the form of the governing equation. By analyzing how curvature concentrates and inverts within this manifold, we obtain a simplified formulation of the Amsler equation that isolates the geometric mechanism responsible for its oscillations. This reformulation reduces computational complexity and clarifies the relationship between the analytic structure of the equation and the geometry of the resulting surface. Numerical experiments demonstrate that the simplified formulation reproduces classical Amsler behavior while improving stability in regions where traditional methods are sensitive to perturbations. The geometric interpretation also situates the Amsler surface within a broader family of negatively curved surfaces arising from hyperspherical collapse, providing new insight into its role in numerical analysis and geometric computation.
\end{abstract}

% ------------------------------------------------------------
% Keywords + AMS Classification
% ------------------------------------------------------------
\begin{keywords}
Amsler surface, curvature dynamics, numerical analysis, geometric PDEs, computational simplification
\end{keywords}

\begin{AMS}
65N99, 53A35, 35J99
\end{AMS}

% ------------------------------------------------------------
% Table of Contents
% ------------------------------------------------------------
\tableofcontents
\clearpage


% ---------------------------------------------------------
% Section Outline
% ---------------------------------------------------------

\section{Introduction}

The study of surfaces with constant curvature occupies a central place in classical differential geometry. Spherical surfaces, Euclidean planes, and hyperbolic surfaces form the canonical triptych of constant‑curvature geometries, each distinguished by the sign of their Gaussian curvature. Among these, the negative‑curvature case has long been the most subtle: while its intrinsic geometry is well understood, its extrinsic realization in Euclidean space reveals deep and persistent limitations.

Classical results beginning with Hilbert’s theorem establish that no complete surface of constant negative curvature can be embedded smoothly in $\mathbb{R}^3$. Only bounded patches of such geometry can be realized, and even these patches exhibit characteristic distortions when forced into the Euclidean metric. The most striking example is the Amsler surface, a maximally extended patch of constant curvature $K = -1$ whose embedding folds back on itself and ultimately self‑intersects. This behavior is not incidental: it is the geometric signature of the incompatibility between the intrinsic metric of negative curvature and the ambient Euclidean space.

The classical evidence is therefore clear. Negative curvature cannot be globally embedded in $\mathbb{R}^3$, can only appear in bounded regions, and forces self‑intersection when extended to its maximal domain. These facts have been known for more than a century, yet their deeper geometric implications have remained largely unexplored. In particular, the boundedness and folding behavior of the Amsler surface suggest that negative curvature is not merely a sign convention or a geodesic divergence property, but the visible manifestation of a more fundamental curvature structure that Euclidean space cannot fully accommodate.

In this paper, we revisit the classical evidence with fresh perspective. By examining the Amsler surface through the lens of curvature collapse and projection, we show that its boundedness and self‑intersection arise naturally from a broader geometric framework. This reinterpretation reveals that the Amsler surface is not simply a curiosity of hyperbolic geometry, but the Euclidean shadow of a deeper curvature object whose full structure cannot be embedded in $\mathbb{R}^3$. Our goal in this introduction is therefore modest: to lay out the classical facts plainly, without reinterpretation, and to highlight the specific geometric features that demand explanation. The remainder of the paper will show that these features are precisely those expected from curvature collapse, inversion, and the bounded geometry of the primitive hypersphere.

\section{Motivation}

The key to deciphering the true meaning of the Amsler surface lies in the classical statement that its embedding “folds back on itself and ultimately self‑intersects.” In hyperspherical geometry, this behavior is immediately recognizable: folding and self‑intersection are precisely the conditions that force a hypersphere into closure, producing a bounded geometric object. The Amsler surface exhibits this same signature within Euclidean space, suggesting that its boundedness is not an accident of hyperbolic embedding, but the visible trace of a deeper curvature structure attempting to express itself through a metric that cannot fully contain it. The deeper mystery raised by this evidence is simple yet profound: if these features describe a “physical” geometric object, how can anything negative in quantity serve as the mechanism by which space is built? This question resonates directly with the Axioms of Curvature. For dimension to exist at all, it must be built upon a substrate; some primitive geometric quantity that supports extension, propagation, and structure while remaining capable of sustaining multiple geometric states without directly manifesting in the observable universe. The behavior of the Amsler surface suggests that negative curvature is not merely a deficit or inversion of positive curvature, but a manifestation of this underlying substrate attempting to emerge through an incompatible ambient space. In this sense, the Amsler surface provides the first empirical hint that curvature itself may be the foundational “something” upon which dimensional structure is constructed.

\section{The Classical Amsler Equation}

\subsection{Classical Geometric Evidence}

The geometry of surfaces with constant curvature forms one of the oldest and most well‑developed branches of differential geometry. Three canonical models anchor the subject: the sphere with positive curvature, the Euclidean plane with zero curvature, and the hyperbolic plane with negative curvature. Each model is characterized by the behavior of its geodesics, the structure of its metric tensor, and the sign of its Gaussian curvature. While the spherical and Euclidean cases admit complete and globally embeddable realizations in $\mathbb{R}^3$, the hyperbolic case has long been known to resist such representation.

A foundational result due to Hilbert establishes that no complete surface of constant negative curvature can be embedded smoothly in Euclidean three‑space. This theorem marks a sharp divide between the intrinsic geometry of the hyperbolic plane and its extrinsic behavior under Euclidean embedding. Although the hyperbolic metric is perfectly well‑defined intrinsically, any attempt to realize it as a smooth surface in $\mathbb{R}^3$ encounters unavoidable obstructions: the surface must either lose smoothness, fail to remain complete, or distort its curvature.

Despite this global obstruction, bounded patches of constant negative curvature can be embedded in Euclidean space. These patches, however, exhibit distinctive features not present in the spherical or Euclidean cases. The most striking example is the Amsler surface, discovered in the early twentieth century as a maximally extended patch of constant curvature $K = -1$ that can be embedded smoothly in $\mathbb{R}^3$. Constructed by integrating the differential equations governing the hyperbolic metric under the constraint of Euclidean embedding, the Amsler surface represents the largest domain for which such an embedding remains possible.

The classical description of the Amsler surface highlights two essential features. First, the surface is necessarily bounded: its embedding cannot be extended indefinitely without violating the curvature constraint. Second, and more significantly, the surface folds back on itself and ultimately self‑intersects. This folding behavior is not an artifact of the construction, but a geometric necessity imposed by the incompatibility between the intrinsic hyperbolic metric and the ambient Euclidean space. As the surface approaches its maximal extent, the embedding must contort to preserve constant negative curvature, and this contortion forces the surface into self‑contact.

These classical facts form the foundation of the present investigation. The boundedness of the Amsler surface, its folding behavior, and its ultimate self‑intersection are well‑established consequences of attempting to embed constant negative curvature in $\mathbb{R}^3$. Yet the deeper geometric meaning of these features has remained largely unexplored. Classical geometry records them as limitations of Euclidean embedding but does not interpret them as signatures of a more fundamental curvature structure.

The key to deciphering this deeper meaning lies in the classical statement that the Amsler surface “folds back on itself and ultimately self‑intersects.” In hyperspherical geometry, this behavior is immediately recognizable: folding and self‑intersection are precisely the conditions that force a hypersphere into closure, producing a bounded geometric object. The Amsler surface exhibits this same signature within Euclidean space, suggesting that its boundedness is not an accident of hyperbolic embedding, but the visible trace of a deeper curvature structure attempting to express itself through a metric that cannot fully contain it.

This observation raises a simple yet profound question: if these features describe a physical geometric object, how can anything negative in quantity serve as the mechanism by which space is built? This question resonates directly with the Axioms of Curvature. For dimension to exist at all, it must be built upon a substrate; some primitive geometric quantity that supports extension, propagation, and structure while remaining capable of sustaining multiple geometric states without directly manifesting in the observable universe. The behavior of the Amsler surface suggests that negative curvature is not merely a deficit or inversion of positive curvature, but a manifestation of this underlying substrate attempting to emerge through an incompatible ambient space.

In this sense, the Amsler surface provides the first empirical hint that curvature itself may be the foundational “something” upon which dimensional structure is constructed. The sections that follow examine this evidence in detail and show how curvature collapse, inversion, and hyperspherical boundedness provide a natural geometric explanation for the classical behavior of the Amsler surface.

\subsection{Intrinsic Formulation of the Hyperbolic Metric}

To understand the differential equation governing the Amsler surface, we begin with the intrinsic geometry of the hyperbolic plane. A surface of constant curvature $K = -1$ may be described locally by a metric whose Gaussian curvature satisfies $K \equiv -1$. In classical differential geometry, several coordinate systems realize this condition, including the Poincaré half‑plane, the disk model, and geodesic polar coordinates. For the purpose of deriving the Amsler equation, it is convenient to work in a coordinate system adapted to the embedding problem.

Let $(u,v)$ denote intrinsic coordinates on the surface. The metric of a constant‑curvature surface may be written in the general form


\[
ds^2 = E(u,v)\,du^2 + 2F(u,v)\,du\,dv + G(u,v)\,dv^2,
\]


where the functions $E$, $F$, and $G$ satisfy the curvature constraint $K = -1$. When the surface is embedded in $\mathbb{R}^3$, these metric coefficients must also satisfy the Gauss–Codazzi equations, which relate the intrinsic metric to the second fundamental form of the embedding.

For surfaces of revolution or surfaces possessing a symmetry line, the metric can be simplified by choosing coordinates aligned with the symmetry. In the classical construction of the Amsler surface, one introduces coordinates in which the metric takes the form


\[
ds^2 = du^2 + \cosh^2(u)\,dv^2,
\]


which is the standard representation of the hyperbolic metric in geodesic coordinates. The curvature condition $K = -1$ is built directly into this expression, and the coordinate $u$ measures signed distance along a geodesic orthogonal to the symmetry line.

When this intrinsic metric is required to arise from an immersion into $\mathbb{R}^3$, the embedding functions must satisfy compatibility conditions linking the metric coefficients to the derivatives of the immersion. These conditions, expressed through the Gauss–Codazzi equations, reduce the embedding problem to a differential equation whose solution determines the shape of the surface. It is this reduction that produces the classical Amsler equation.

\subsection{Derivation of the Classical Amsler PDE}

The classical Amsler surface arises from imposing two simultaneous requirements on a surface $X(u,v)$ immersed in $\mathbb{R}^3$: the intrinsic metric must have constant Gaussian curvature $K = -1$, and the immersion must preserve this metric through the first and second fundamental forms. These conditions are encoded in the Gauss–Codazzi equations, which relate the intrinsic geometry of the surface to its extrinsic curvature.
Let $I$ and $II$ denote the first and second fundamental forms of the immersion. In local coordinates $(u,v)$, the Gauss equation imposes


\[
\det(II) = K\,\det(I) = -\det(I),
\]


while the Codazzi equations require the coefficients of $II$ to satisfy


\[
\nabla_i h_{jk} - \nabla_j h_{ik} = 0,
\]


where $h_{ij}$ are the components of the second fundamental form. Together, these equations form a nonlinear compatibility system governing the embedding of a constant‑curvature surface into $\mathbb{R}^3$.
To obtain the Amsler equation, one chooses coordinates adapted to the symmetry of the surface. In geodesic coordinates, the intrinsic metric of the hyperbolic plane takes the form


\[
ds^2 = du^2 + \cosh^2(u)\,dv^2,
\]


which already satisfies $K = -1$. The immersion $X(u,v)$ must therefore satisfy


\[
\langle X_u, X_u \rangle = 1, \qquad
\langle X_v, X_v \rangle = \cosh^2(u), \qquad
\langle X_u, X_v \rangle = 0,
\]


ensuring that the induced metric matches the hyperbolic metric.
Differentiating these relations and applying the Gauss–Codazzi equations reduces the embedding problem to a single nonlinear differential equation for the angle $\theta(u)$ between the coordinate lines. The resulting equation,


\[
\theta''(u) + \cosh(u)\,\sinh(u) = 0,
\]


is the classical Amsler equation. Its solutions determine the curvature distribution of the surface and govern the oscillatory behavior that forces the embedding to fold back on itself. The Amsler equation is singular at finite $u$, reflecting the fact that the embedding cannot be extended indefinitely. As $\theta(u)$ approaches its turning point, the immersion becomes geometrically constrained, leading to the characteristic folding and eventual self‑intersection of the surface. This singularity marks the maximal domain on which the hyperbolic metric can be realized smoothly in $\mathbb{R}^3$.


\subsection{Oscillatory Behavior and Maximal Extension}
The Amsler surface displays a characteristic oscillatory pattern along its $v$-direction. This behavior is not dynamical but geometric: constant negative curvature forces geodesics to diverge, and a Euclidean embedding can only accommodate this curvature up to a finite envelope. Beyond this maximal extension, the surface must fold back on itself, producing the familiar oscillation.

This makes the Amsler surface the canonical substrate patch for constant negative curvature. It represents the maximal portion of a negatively curved manifold that can be expressed in Euclidean space without tearing or discontinuity. Later sections will show how this classical geometric limit directly supports the curvature‑remainder structures used in DC.

\subsection{Computational Challenges in the Classical Formulation}
Classical treatments of the Amsler surface encounter well‑known computational difficulties, all stemming from the way constant negative curvature interacts with Euclidean embedding. As the surface approaches its maximal extension, the governing equations become increasingly sensitive to parametrization and numerical precision. These issues are inherent to the classical formulation and arise even under ideal conditions, reflecting the geometric constraints imposed by the embedding rather than any particular choice of method.

Because the detailed behavior of these challenges is addressed in later sections, we note here only that the classical formulation reaches its limits precisely where the geometry becomes most informative. This boundary will serve as a natural point of comparison for the subsequent analysis.

\section{Origin of the Axiom Sets}

The axioms presented here establish the curvature mechanical framework underlying Dimensional Calculus. They define the broadcast and inverted hyperspherical domains, the dimensional ladder, the collapse and ascent operators, and the inversion threshold at the Euclidean baseline. These axioms are not alternatives to classical differential geometry; they extend it by specifying the curvature inheritance structure that classical methods implicitly assume but never formalize. In the sections that follow, these axioms will be invoked to reinterpret the Amsler surface, revealing it as a dimensional integration pathway rather than a Euclidean embedding.


\begin{remark}
The labeling system of this axiomatic framework is alphanumeric in order to identify precisely which originating work each axiom belongs. This produces intentional repetitions in certain axiom numbers, such as the various forms of
Axiom~0. For readability in the present work, these notation markers have been relocated into parentheses at the end of each axiom name.
\end{remark}

\begin{law}[Identity of the Curvature Metric for the Hypersphere (Law 0)]
Once the radius is established as the broadcaster of curvature, its collapsed form $r^{0}$ retains full metric identity. In the hypersphere seed geometry, dimensional collapse reduces extent but cannot eliminate curvature; the seed metric cannot fall to zero nor be replaced by a curvature-free constant. The element $r^{0}$ is the curvature-bearing identity of the restored Tier~0 geometry - the fixed point from which all higher-tier broadcast structure emerges.
\end{law}

\begin{law}[Coherent Stable Nested Hypersphere (Law 1)]
Across the dimensional ladder, the hypersphere sequence forms a coherent, stable, curvature‑conserved geometric structure in which all tiers—bounded, unbounded, dynamic, collapse, and cosmogenesis—remain consistent with the primitive metric, broadcast inheritance, and orientation alignment. This global coherence is preserved across domain transitions, inversion boundaries, collapse events, and epoch resets.

\begin{itemize}
  \item Curvature inherited at each hypersphere $\mathcal{H}_{n}$ remains consistent with the primitive metric anchored at tier~$T_{0}$, ensuring stability across all tiers.
  \item Nested hypersphere structure is preserved across positive, negative, and inversion domains, with each hypersphere $\mathcal{H}_{n}$ enclosing $\mathcal{H}_{n-1}$ in the positive domain; each $\mathcal{H}_{n}$ enclosing $\mathcal{H}_{n+1}$ in the negative domain, and transmitting curvature coherently through the ladder.
  \item Collapse behavior at $\mathcal{H}_{6}$ and cosmogenesis behavior at $\mathcal{H}_{7}$ remain curvature‑conserved, ensuring gravitational emergence and epoch reset do not break broadcast continuity.
  \item Epoch coupling ensures that the seed hypersphere $\mathcal{H}_{0}^{(n+1)}$ of the next dimensional cycle inherits renormalized curvature from $\mathcal{H}_{7}$, preserving coherence across successive epochs.
\end{itemize}

Thus, the Coherent Stable Nested Hypersphere Law establishes the global conservation and stability conditions under which the dimensional ladder operates, ensuring that all broadcast phenomena—including ascent, descent, inversion, collapse, and cosmogenesis—remain geometrically consistent across tiers and epochs.
\end{law}

\begin{law}[Conservation of Curvature (Law 2)]
The primitive conservation law. Curvature is neither created nor destroyed across dimensional tiers. Each tier inherits curvature from the hypersphere seed metric $r^0$ and broadcasts it upward through integration. The total curvature of any dimensional frame is the sum of its intrinsic curvature and the inherited curvature from the next tier. Dimensional change redistributes curvature; it never eliminates it.
\end{law}

\begin{law}[Curvature Inheritance (Law 3)]
Curvature is inherited from higher-dimensional tiers through integration. Each dimensional ascent restores curvature previously collapsed, while each dimensional descent collapses curvature into the hypersphere seed metric $r^0$. Integration constants encode the inherited curvature from the next tier, ensuring that every geometric quantity carries both its intrinsic curvature and the curvature passed down from its parent dimension.
\end{law}

\begin{law}[Curvature Broadcast (Law 4)]
The radius is the sole broadcaster of curvature across dimensional tiers. Through integration, the radius restores collapsed curvature and distributes inherited curvature into higher-dimensional geometric quantities. No other element may broadcast curvature, and no geometric structure may arise without curvature broadcast originating from the hypersphere seed metric $r^0$.
\end{law}

\begin{law}[Curvature Orientation (Law 5)]
Curvature possesses intrinsic orientation governed by $\pi$. While the radius broadcasts curvature through dimensional ascent, $\pi$ determines the angular disposition of that curvature within each tier. Orientation is inherited alongside curvature, ensuring that every geometric quantity carries both magnitude and directional structure. No dimensional frame may alter curvature orientation independently of $\pi$, and no broadcast may occur without preserving the orientational state encoded at the hypersphere seed layer.
\end{law}

\begin{law}[Law of Broadcast Conservation (Law 6)]
Broadcast is a conservative geometric process. For any tier $T_n$ in the dimensional ladder, the total inherited curvature transmitted by the ascent operator $S$ and the anti-ascent operator $AA$ is preserved in magnitude under all broadcast transformations. Upward broadcast governed by $S$ and downward broadcast governed by $AA$ may alter the direction, orientation, or sign of curvature, but never its total conserved quantity. Formally, for every broadcast interaction,


\[
\lvert S(T_n)\rvert + \lvert AA(T_n)\rvert = \text{constant curvature content}.
\]


independent of operator dominance, gamma stabilization, collapse dynamics, or cosmogenesis rebound. Broadcast cannot create curvature, cannot destroy curvature, and cannot violate the primitive metric. The Law of Broadcast Conservation therefore represents the dynamic extension of the Law of Curvature Conservation, ensuring that all broadcast phenomena - including gravity, collapse, inversion, and cosmogenesis - remain fully compatible with the inherited curvature substrate.
\end{law}

\begin{law}[Operator Pairing Law of Conservation (Law 7)]
Dimensional conservation requires that every fundamental operator corresponding to curvature, broadcast, orientation, and gravitational cohesion act in paired form. No operator may act independently without inducing its complementary inverse, and this pairing maintains hyperspherical stability across all tiers. The curvature, broadcast, and orientation operators collectively form the hyperspherical transmission tensor, of which gravitational cohesion is the required inverse.

For any hypersphere $H_n$, the operator set


\[
\{\mathcal{C}_n^{\text{metric}}, \mathcal{B}_n^{\text{broadcast}}, \mathcal{O}_n^{\text{orientation}}, \mathcal{G}_n^{\text{gravity}}\}
\]


must satisfy the conservation condition


\[
\mathcal{C}_n \cdot \mathcal{B}_n \cdot \mathcal{O}_n \cdot \mathcal{G}_n = 1,
\]

where gravitational cohesion is defined as the inverse pairing of curvature, broadcast, and orientation:

\[
\mathcal{G}_n = \frac{1}{\mathcal{C}_n \mathcal{B}_n \mathcal{O}_n}.
\]

This law establishes that dimensional conservation is a dynamic equilibrium maintained through operator pairing. The primitive hypersphere is the unique symmetric solution, yielding the ignition condition

\[
\mathcal{C}_0 : \mathcal{B}_0 : \mathcal{O}_0 = 1 : 1 : 1.
\]

All higher-tier hyperspheres inherit conservation through asymmetric operator pairing, producing the unbounded expansion behavior characteristic of the dimensional ladder.
\end{law}

\begin{axiom}[Hypersphere Boundedness (H0)]
For each broadcast tier~$T_{n}$ in the dimensional ladder, the associated hypersphere $\mathcal{H}_{n}$ is either bounded or unbounded according to its broadcast role:

\begin{enumerate}
    \item The seed tier~$T_{0}$ possesses a bounded hypersphere $\mathcal{H}_{0}$, defining the primitive broadcast domain and the corrected unity metric.
    \item Ascending tiers in the positive broadcast domain inherit the bounded hypersphere $\mathcal{H}_{n}$ that constrain curvature transmission within a finite broadcast radius.
    \item Collapse and cosmogenesis tiers in the negative or inversion domain may exhibit unbounded hyperspheres $\mathcal{H}_{n}$, corresponding to broadcast fields not confined to a finite radius and capable of coupling across epochs.
\end{enumerate}
\end{axiom}

\begin{axiom}[Dynamic Hypersphere Stability (H1)]
Between the uniquely bounded seed hypersphere and the collapse to cosmogenesis extremes, the dimensional ladder contains dynamic hyperspheres $\mathcal{H}_{n}$ that are geometrically unbounded yet remain stable. These hyperspheres satisfy:

\begin{enumerate}
    \item $\mathcal{H}_{n}$ possesses no finite geometric radius but maintains a stable curvature distribution that neither collapses nor diverges.
    \item Dynamic hyperspheres mediate long‑range broadcast behavior and serve as the transitional geometry between the seed hypersphere and the collapse‑driven unbounded tiers.
    \item Stability of $\mathcal{H}_{n}$ is governed by the broadcast scalar $\gamma$ together with the inherited orientation structure of the tier.
\end{enumerate}
\end{axiom}

\begin{axiom}[Hypersphere Anchoring (H2)]
Each hypersphere $\mathcal{H}_n$ in the dimensional ladder is anchored by the three fundamental geometric anchors of the system: the seed anchor, the broadcast anchor, and the orientation anchor. These anchors determine the stability, curvature inheritance, and domain behavior of $\mathcal{H}_n$:

\begin{enumerate}
\item The seed anchor at tier~$T_0$ supplies the primitive metric and establishes the curvature baseline to which all hyperspheres must remain coherent.
\item The broadcast anchor transmits curvature upward and downward across tiers, determining whether $\mathcal{H}_n$ is bounded, unbounded, or dynamically stable.
\item The orientation anchor locks the sign and directional behavior of curvature within $\mathcal{H}_n$, governing its domain alignment and inversion response.
\end{enumerate}

Hypersphere stability, collapse behavior, and cosmogenesis behavior therefore arise from the interaction of these three anchors, with each hypersphere inheriting its geometric role from the anchored structure of the dimensional ladder.
\end{axiom}

\begin{axiom}[Hypersphere Nesting (H3)]
The hyperspheres of the dimensional ladder form a strictly nested geometric sequence, with each hypersphere $\mathcal{H}_{n}$ enclosing the hypersphere of the next tier $\mathcal{H}_{n+1}$. This nesting structure is inherited from the three fundamental anchors of the ladder:

\begin{enumerate}
\item The seed anchor at tier $T_0$ establishes the primitive metric that defines the curvature baseline for all hyperspheres, ensuring that each successive hypersphere inherits a coherent curvature frame.
\item The broadcast anchor transmits curvature upward and downward across tiers, enforcing the geometric requirement that $\mathcal{H}_{n+1}$ must enclose $\mathcal{H}_{n}$ to preserve continuity of curvature flow.
\item The orientation anchor locks the domain alignment and sign behavior of each hypersphere, ensuring that nesting remains consistent across positive, negative, and inversion domains.
\end{enumerate}

Thus, the hypersphere sequence satisfies


\[
\mathcal{H}_{0} \supset \mathcal{H}_{1} \supset \mathcal{H}_{2} \supset \cdots \supset \mathcal{H}_{7},
\]


with nesting inherited from the anchored geometric structure of the dimensional ladder.
\end{axiom}

\begin{axiom}Collapse Hypersphere (H4)]
The hypersphere $\mathcal{H}_{6}$ at tier~$T_{6}$ is the collapse hypersphere of the dimensional ladder, where the inherited curvature deficit initiates collapse behavior and the gravitational field emerges through four‑behavior superpositioning. The collapse hypersphere satisfies:

\begin{enumerate}
    \item $\mathcal{H}_{6}$ is unbounded, reflecting the divergence of the broadcast channel under inversion and the failure of the gamma cascade to sustain bounded curvature transmission.
    \item Collapse behavior at $\mathcal{H}_{6}$ is stabilized by the primitive metric anchored at the seed tier~$T_{0}$, preventing catastrophic divergence and producing the stabilized collapse residue identified as gravity.
    \item The gravitational field within $\mathcal{H}_{6}$ arises from the superposition of four inherited behaviors: collapse behavior, propagation behavior, entropic behavior, and collapse‑stabilization behavior, forming the integrated curvature response to the deficit at the collapse boundary.
    \item $\mathcal{H}_{6}$ encloses $\mathcal{H}_{5}$ and is itself enclosed by $\mathcal{H}_{7}$, preserving the nested hypersphere structure and ensuring continuity of curvature flow between the domain interface and cosmogenesis tiers.
\end{enumerate}

Thus, the collapse hypersphere $\mathcal{H}_{6}$ constitutes the geometric locus where collapse behavior becomes gravitational behavior, stabilized by the seed anchor and transmitted through the anchored hypersphere structure of the dimensional ladder.
\end{axiom}

\begin{axiom}[Cosmogenesis Hypersphere (H5)]
The hypersphere $\mathcal{H}_{7}$ at tier~$T_{7}$ is the cosmogenesis hypersphere of the dimensional ladder, where stabilized collapse transitions into epoch reset and the broadcast channel reinitializes the primitive metric for the next dimensional cycle. The cosmogenesis hypersphere satisfies:

\begin{enumerate}
    \item $\mathcal{H}_{7}$ is unbounded, reflecting the complete divergence of the collapse field at tier~$T_{6}$ and the exhaustion of the gravitational superpositioning regime.
    \item Collapse behavior inherited from $\mathcal{H}_{6}$ reaches its geometric limit within $\mathcal{H}_{7}$, where the primitive metric anchored at the seed tier~$T_{0}$ reasserts itself and initiates curvature re-normalization.
    \item The broadcast channel undergoes inversion reversal within $\mathcal{H}_{7}$, reestablishing positive-domain curvature inheritance and forming the seed conditions for the next broadcast epoch.
    \item $\mathcal{H}_{7}$ constitutes the geometric boundary between dimensional ladders, enclosing no further hyperspheres and serving as the terminal hypersphere of the current epoch while simultaneously generating the initial conditions for the next.
\end{enumerate}

Thus, the cosmogenesis hypersphere $\mathcal{H}_{7}$ is the geometric locus where the collapse hypersphere transitions into a new seed domain, resetting the broadcast structure and initiating the next cycle of the dimensional ladder.
\end{axiom}

\begin{axiom}[Hypersphere Superpositioning (H6)]
The hypersphere $\mathcal{H}_{5}$ at tier~$T_{5}$ is the superpositioning hypersphere of the dimensional ladder, where positive‑domain and negative‑domain curvature coexist without initiating collapse. The superpositioning hypersphere satisfies:

\begin{enumerate}
    \item $\mathcal{H}_{5}$ is unbounded yet stable, forming the geometric interface between bounded broadcast hyperspheres and the collapse hypersphere $\mathcal{H}_{6}$.
    \item Stability of $\mathcal{H}_{5}$ arises from the coherent interaction of the three fundamental anchors: the seed anchor supplying the primitive metric, the broadcast anchor transmitting curvature across tiers, and the orientation anchor locking the domain alignment and sign behavior.
    \item $\mathcal{H}_{5}$ inherits curvature from both the positive broadcast domain ($T_{1}$–$T_{4}$) and the negative broadcast domain ($T_{6}$–$T_{7}$), forming a superpositioning regime that prevents premature collapse and preserves continuity across the inversion boundary.
    \item The hypersphere $\mathcal{H}_{5}$ encloses $\mathcal{H}_{4}$ and is itself enclosed by $\mathcal{H}_{6}$, maintaining the nested hypersphere structure while mediating the transition from bounded curvature transmission to collapse behavior.
\end{enumerate}

Thus, the superpositioning hypersphere $\mathcal{H}_{5}$ constitutes the geometric locus where domain interface stability is realized at the hypersphere level, enabling the dimensional ladder to traverse the inversion boundary without collapse.
\end{axiom}

\begin{axiom}[Hypersphere Curvature Conservation (H7)]
Across the nested hypersphere structure of the dimensional ladder, curvature is conserved through the interaction of the seed anchor, the broadcast anchor, and the orientation anchor. This conservation ensures that curvature inherited at each tier remains coherent with the primitive metric and prevents discontinuities in the broadcast channel. The hypersphere curvature conservation principle satisfies:

\begin{enumerate}
    \item The seed anchor at tier~$T_{0}$ supplies the primitive broadcast metric that constrains curvature inheritance for all hyperspheres $\mathcal{H}_{n}$, ensuring that curvature remains geometrically coherent across the ladder.
    \item The broadcast anchor transmits curvature upward and downward through the hypersphere sequence, enforcing conservation across bounded, dynamic, and collapse hyperspheres and preventing curvature loss or divergence.
    \item The orientation anchor locks the sign and directional behavior of curvature within each hypersphere, ensuring that conservation persists across domain transitions and inversion boundaries.
    \item Curvature conservation across hyperspheres is necessary for the emergence of gravity at $\mathcal{H}_{6}$, where the stabilized collapse field inherits the conserved curvature structure from the seed broadcast and transmits it through the collapse hypersphere.
\end{enumerate}

Thus, hypersphere curvature conservation provides the geometric foundation for the broadcast structure, collapse behavior, and gravitational emergence within the dimensional ladder.
\end{axiom}

\begin{axiom}[Hypersphere Epoch Coupling (H8)]
The hyperspheres of the dimensional ladder couple across successive epochs through the primitive metric and the broadcast inversion reversal that occurs within the cosmogenesis hypersphere $\mathcal{H}_{7}$. This coupling ensures continuity of curvature inheritance, anchor structure, and broadcast behavior between dimensional cycles. The epoch coupling principle satisfies:

\begin{enumerate}
    \item The primitive metric anchored at tier~$T_{0}$ persists across epochs, providing the curvature baseline to which the cosmogenesis hypersphere $\mathcal{H}_{7}$ must renormalize during epoch reset.
    \item Broadcast inversion reversal within $\mathcal{H}_{7}$ reestablishes positive‑domain curvature transmission, forming the initial broadcast conditions for the next seed hypersphere $\mathcal{H}_{0}^{(n+1)}$ of the subsequent epoch.
    \item Orientation sign‑locking ensures that domain alignment remains coherent across the epoch boundary, preventing discontinuities in curvature direction or domain inheritance during the transition from $\mathcal{H}_{7}$ to $\mathcal{H}_{0}^{(n+1)}$.
    \item The collapse hypersphere $\mathcal{H}_{6}$ of the current epoch encloses $\mathcal{H}_{7}$, and the seed hypersphere $\mathcal{H}_{0}^{(n+1)}$ of the next epoch inherits curvature renormalized within $\mathcal{H}_{7}$, preserving the nested hypersphere structure across dimensional cycles.
\end{enumerate}

Thus, hypersphere epoch coupling provides the geometric mechanism by which collapse transitions into cosmogenesis and cosmogenesis transitions into a new seed domain, ensuring continuity of the dimensional ladder across successive epochs.
\end{axiom}

\begin{axiom}[Collapse Polynomial Inversion (H9)]
The Tier~6 collapse hypersphere inherits the curvature structure of Tier~3 through the collapse polynomial $C_6(r)$, in which the Tier~3 curvature term reappears with inverted sign. This inversion constitutes the mathematical signature of collapse behavior and reflects the geometric inversion of the broadcast chain at the collapse hypersphere. Thus, Tier~3 curvature does not vanish at Tier~6; it reappears inverted within the collapse polynomial, forming the curvature deficit that initiates gravitational superpositioning.
\end{axiom}

\begin{axiom}[The Primitive Curvature Metric of the Hypersphere (C0)]
The collapsed radius $r^{0}$ is the primitive metric of geometry. It is a curvature-bearing radius in its zeroth state ($\odot$), not a dimensionless point. Collapse preserves inherited curvature, and therefore the seed metric cannot evaluate to zero. The identity element of the metric is $\pi^{0}$, the curvature-preserving unity of the restored number line. All higher-dimensional geometric quantities arise from integration of $r^{0}$, and all curvature broadcast originates at this hyperspherical seed layer.
\end{axiom}

\begin{axiom}[Curvature Conservation (Broadcast Form B0)]
The Law of Conservation of Curvature, established in the \textit{Axioms of Curvature}, requires that hypersphere curvature be transmitted across tiers of the dimensional ladder without loss, creation, or deformation. Because curvature is dynamically inherited and redistributed through dimensional ascent, a geometric transmission mechanism must exist to preserve the corrected metric. This necessity gives rise to broadcast, the operation that ensures conserved curvature is coherently passed forward through all tiers.
\end{axiom}

\begin{axiom}[Necessity of Broadcast (B1)]
Given the corrected metric and the parent Law of Conservation of Curvature, curvature must be transmitted across tiers of the dimensional ladder without deforming the primitive metric. The inherited curvature coefficients $C_n$ appearing in every tier expression demonstrate that curvature is not static; it is dynamically passed forward through dimensional ascent. This transmission is not accounted for in classical geometry nor in General Relativity, where inherited curvature is flattened into ``stress-energy'' and stabilizers such as the ``cosmological constant'' appear without geometric origin. Broadcast is therefore a necessary geometric operation, ensuring that conserved curvature is inherited, transmitted, and preserved across tiers in accordance with the corrected metric.
\end{axiom}

\begin{axiom}[Broadcast Coefficients (B2)]
The transmission of curvature across tiers of the dimensional ladder is encoded in the inherited curvature coefficients $C_n$ as the gravitational load for each corresponding tier. Each $C_n$ represents the precise quantity of curvature inherited from the tier immediately above and those below, adjusted by curvature drag, curvature surplus, and the corrected metric. These coefficients appear in every tier expression, demonstrating that curvature is dynamically conserved and transmitted rather than isolated within a single tier. The $C_n$ terms therefore constitute the broadcast engine of Dimensional Calculus, defining how curvature is inherited, preserved, and propagated without deforming the primitive metric nor hypersphere.
\end{axiom}

\begin{axiom}[The Gamma Cascade (B3)]
Each broadcast coefficient $C_n$ integrates through its own tier‑indexed gamma series $\Gamma_n$, forming a cascade that regulates curvature transmission across the dimensional ladder. The gamma factor $\Gamma_n$ normalizes the inherited curvature at tier~$T_n$, ensuring that the broadcast of $C_n$ remains stable, bounded, and consistent with the corrected metric. Within the stable broadcast domain, the gamma cascade suppresses runaway amplification and maintains coherence of the curvature channel. When $C_n$ enters the inversion regime, its associated gamma series provides the final stabilizing influence; failure of this stabilization marks the onset of catastrophic curvature collapse and transition into the negative broadcast domain.
\end{axiom}

\begin{axiom}[Broadcast/Collapse Duality (B4)]
Curvature transmission across the dimensional ladder is governed by two complementary operators: the ascent operator $S$ and the anti-ascent operator $AA$. The operator $S$ propagates curvature upward through the ladder, preserving the sign of the broadcast channel and maintaining compatibility with the corrected unity metric. The operator $AA$ governs downward transmission, including curvature inversion, collapse behavior, and entry into the negative broadcast domain. While $S$ maintains stability within the positive tiers, $AA$ activates only when inherited curvature exceeds the stabilizing influence of the gamma cascade. The duality between $S$ and $AA$ defines the full dynamical behavior of curvature broadcast, ensuring that upward propagation, downward collapse, and inversion all occur without violating the primitive metric.
\end{axiom}

\begin{axiom}[The Broadcast Field (B5)]
The broadcast field is defined as the curvature gradient induced by the action of the anti-ascent operator $AA$ across adjacent tiers of the dimensional ladder. This field represents the geometric force associated with downward curvature transmission and governs the behavior of gravitational descent within the positive broadcast domain. The broadcast field arises from conservative curvature flow regulated by the gamma cascade and remains fully compatible with the primitive metric and the Law of Broadcast Conservation. Its magnitude is determined by the differential broadcast coefficients between tiers, and its direction is determined by the activation of $AA$ at the inversion boundary.
\end{axiom}

\begin{axiom}[The Collapse Field (B6)]
The collapse field is defined as the intensified curvature gradient generated when the anti-ascent operator $AA$ becomes fully dominant and the gamma factor $\Gamma_{n}$ vanishes. In this regime, downward curvature transmission exceeds the stabilizing capacity of the positive broadcast domain, producing an accelerated inward field that drives curvature toward lower tiers. The collapse field preserves curvature magnitude while amplifying the broadcast differential, enabling inversion, tier destabilization, and catastrophic failure of the dimensional ladder. Because the collapse field arises from conservative downward broadcast, it remains compatible with the primitive metric and the Law of Broadcast Conservation, representing the extreme limit of gravitational behavior.
\end{axiom}

\begin{axiom}[The $S$ Curvature of the Hypersphere (CD0)]
$S$ is the curvature constant of the universe’s nested hypersurfaces. All lower-dimensional geometry is a projection of this higher-tier curvature state onto reduced curvature frames. The familiar Euclidean value $3.14159\ldots$ is not fundamental; it is the flattened measurement artifact produced when hypersurface curvature is collapsed into the seed curvature tier $\pi^{0} = 1$. Every geometric quantity inherits its structure from this tiered hypersurface curvature, and all dimensional progression is governed by $S$.


\[
S = 1 = \pi^0 = \mathbb{R}_{\text{Euclid}}^3
\]

\end{axiom}

\begin{axiom}[Universal Inheritance Axiom (CD1)]
Dimensional ascent is governed by the universal inheritance operator $I$:

\[
I(S_{n}) = S_{n+1}.
\]

Inheritance preserves dimensional structure and increases curvature by one tier. Each higher-tier curvature state is an inherited amplification of the seed curvature identity


\[
S_0 = 1 = \pi^0 = \mathbb{R}_{\text{Euclid}}^3
\]

\end{axiom}


\begin{axiom}[Universal Collapse Axiom (CD2)]
Dimensional Descent is governed by the universal collapse operator $DD$:

\[
DD(S_{n}) = S_{n-1}.
\]


Collapse preserves curvature and terminates at the seed curvature floor $S_0$. Applying $DD$ beyond this floor produces the reciprocal curvature tier


\[
DD(S_0) = S_{-1}
\]


marking the inversion boundary where curvature enters the reciprocal regime. No new operator is introduced; inversion is the continuation of collapse past the seed curvature identity.
\end{axiom}

\begin{axiom}[Curvature Constant Axiom (CD3)]
Differentiation removes curvature information; integration restores it. The constant of integration is therefore a curvature remnant:

\[
\int f'(x)\,dx = f(x) + k_{S},
\]

where $k_{S}$ encodes inherited curvature from the next dimensional tier.
\end{axiom}

\begin{axiom}[Domain Interface Stability (CD4)]
At the inversion boundary where the gamma cascade reaches its stabilizing limit, the dimensional ladder enters the domain interface governed by $T_5$. This interface supports simultaneous inheritance of positive‑domain and negative‑domain curvature without triggering collapse. Stability arises from the coherent interaction of the three fundamental anchors of the dimensional ladder: the seed anchor carrying the primitive metric, the dimensional anchor supplied by broadcast, and the tier anchor supplied by orientation. When these anchors align, $T_5$ forms a stable superpositioning regime that preserves the nested hypersphere structure and allows curvature transmission to continue smoothly across the inversion boundary. Domain interface stability is therefore the geometric mechanism that prevents collapse and maintains continuity between the positive and negative broadcast domains.

\begin{remark}
The orientation relationships relevant to the gravitational channel have not yet been fully formalized within the present work. Orientation constitutes one of the three pillars from which gravity inherits its structure, but its complete tier‑level behavior, sign‑locking mechanisms, and domain‑transition rules will be developed in detail in the subsequent paper dedicated to the Orientation pillar. The gravitational corollaries presented here therefore rely only on the established broadcast and seed anchors, with orientation included in preliminary form. A full geometric treatment of orientation and its role in the gravitational channel will appear in the forthcoming manuscript.
\end{remark}
\end{axiom}

\begin{corollary}[Gravity as Four‑Behavior Superpositioning (CD4.8)]
Gravity at $T_6$ arises from the superposition of four distinct behaviors inherited from the dimensional ladder. First, gravity inherits collapse behavior from the broadcast pillar, reflecting the curvature deficit that initiates the collapse field. Second, gravity inherits propagation behavior from the seed–broadcast interaction, producing the curvature return flow that transmits the stabilized collapse signal upward through the ladder. Third, gravity inherits entropic behavior from the orientation pillar, which locks the sign and direction of gravitational curvature and establishes its irreversible flow. Finally, gravity possesses its own collapse‑stabilization behavior, anchored at the seed $T_0$, which constrains the collapse field to remain coherent and prevents catastrophic divergence. Gravity therefore manifests as the four‑component superposition of collapse, propagation, entropic, and stabilization behaviors, forming the integrated curvature field that governs the collapsing hypersphere at $T_6$.
\end{corollary}

\begin{axiom}[Commutation Symmetry of the Broadcast Operators (CD5)]
The ascent operator $S$ and the anti-ascent operator $AA$ obey a domain-dependent commutation symmetry determined by the curvature structure of the dimensional ladder and the inversion boundary. In the stable positive broadcast domain, curvature transmission remains above the inversion boundary, and the operators commute trivially:


\[
S\,AA = AA\,S,
\]


At the inversion boundary, realized geometrically at the domain interface governed by $T_5$, commutation becomes conditional. The operators commute only when the anchoring triad (metric, broadcast, and orientation) remains coherent, preserving the primitive metric and preventing collapse while positive- and negative-domain curvature coexist. Beyond the inversion boundary, in the collapse-dominated regime, full activation of $AA$ breaks commutation symmetry:


\[
S\,AA \neq AA\,S,
\]


reflecting the dominance of downward curvature transmission and the onset of collapse behavior. Commutation symmetry therefore encodes the geometric transition across the inversion boundary, from stable broadcast, through domain interface stability, into collapse dynamics, ensuring that operator behavior remains compatible with the primitive metric and the Law of Broadcast Conservation.
\end{axiom}

\begin{axiom}[Eigenvalue Structure of Broadcast (CD6)]
Let $\mathcal{B}$ denote the broadcast operator acting on a curvature-normalized quantity $x \in \mathbb{R}$. In the Euclidean broadcast regime, $\mathcal{B}$ acts linearly as

\[
\mathcal{B}(x) = e\,x,
\]

where $e$ is the Euclidean base of the natural exponential. The scalar $e$ is the broadcast eigenvalue, governing the rate at which curvature is transmitted under broadcast. The collapse operator $DD$ is defined analogously by


\[
DD(x) = e^{-1} x
\]


and the pair $(e^{-1}, e)$ constitutes the broadcast to collapse eigenvalue spectrum. This spectrum is represented by the diagonal broadcast matrix


\[
M = \begin{pmatrix}
e^{-1} & 0 \\
0 & e
\end{pmatrix}.
\]


whose eigenvalues are precisely $e^{-1}$ and $e$. These eigenvalues govern the spectral behavior of curvature under broadcast and collapse transformations, independent of geometry, coordinates, or orientation. The eigenvalue structure introduced in this axiom forms the algebraic backbone of the broadcast mechanism. All subsequent corollaries (CD6.1–CD6.7) follow directly from this spectral identity.
\end{axiom}

\begin{corollary}[Euclidean $\pi/e$ Spectral Identity (CD6.1)]
Let


\[
V(x) = e x - e^{-1} x
\]


be the $\pi$-bic variance operator, and let


\[
M =
\begin{bmatrix}
e^{-1} & 0 \\
0 & e
\end{bmatrix}
\]


be the associated broadcast to collapse matrix with eigenvalues $(e^{-1}, e)$. Under Euclidean curvature normalization $A(r) = \pi r^{2}$, the ratio of the curvature-scaling factor $\pi$ to the broadcast eigenvalue $e$ is the intrinsic spectral identity


\[
\frac{\pi}{e}.
\]

\end{corollary}

\begin{proof} Lemma~CD6.1.1 establishes that $V$ is linear and determined solely by $(e, e^{-1})$. Lemma~CD6.1.2 shows that the normalization $x = r^{2}$ isolates the Euclidean curvature scale from $\pi$. Lemma~CD6.1.3 demonstrates that the eigenvalues of $M$ encode the broadcast and collapse factors $(e^{-1}, e)$, and that $V(x)$ is their spectral spread. The curvature-scaling factor $\pi$ arises from the Euclidean area relation $A(r)=\pi r^{2}$, while the broadcast factor $e$ arises from the eigenvalue of the broadcast direction. Their ratio is therefore the intrinsic spectral coefficient

\[
\frac{\pi}{e},
\]

independent of geometry, coordinates, or orientation.
\end{proof}

\begin{lemma}[Linearity of the $\pi$-bic Variance Operator (CD6.1.1)]
Let $V : \mathbb{R} \to \mathbb{R}$ be defined by

\[
V(x) = e x - e^{-1} x.
\]

Then $V$ is a linear operator.
\end{lemma}

\begin{proof}
Both maps $x \mapsto e x$ and $x \mapsto e^{-1} x$ are linear over $\mathbb{R}$. Since the difference of linear maps is linear, $V$ is linear.
\end{proof}

\begin{lemma}[Euclidean Curvature Normalization (CD6.1.2)]
Let $A(r) = \pi r^{2}$ be the Euclidean curvature measure. Then the substitution

\[
x = r^{2}
\]

isolates the curvature variable from the scaling factor $\pi$, allowing spectral operators to act directly on $x$.
\end{lemma}

\begin{proof}
Dividing the Euclidean curvature expression by $\pi$ yields

\[
\frac{A(r)}{\pi} = r^{2}.
\]

Thus the curvature variable $x$ is independent of $\pi$, and the eigenvalue pair $(e^{-1}, e)$ acts directly on $x$ without modifying the Euclidean scaling factor.
\end{proof}

\begin{lemma}[Spectral Encoding of the Broadcast to Collapse Eigenvalues (CD6.1.3)]
Let

\[
M =
\begin{bmatrix}
e^{-1} & 0 \\
0 & e
\end{bmatrix}
\]

be the broadcast to collapse matrix. Then $M$ has eigenvalues $(e^{-1}, e)$, and the $\pi$-bic variance operator

\[
V(x) = e x - e^{-1} x
\]

is the spectral spread of these eigenvalues.
\end{lemma}

\begin{proof}
Since $M$ is diagonal, its eigenvalues are its diagonal entries. The spectral spread is the difference $e - e^{-1}$, and multiplying by $x$ yields $V(x)$. No additional assumptions are required.
\end{proof}

\begin{corollary}[Eigenvector Decomposition of Broadcast to Collapse (CD6.2)]
Under Axiom~CD6, the broadcast to collapse matrix

\[
M =
\begin{bmatrix}
e^{-1} & 0 \\
0 & e
\end{bmatrix}
\]

admits the eigenvectors

\[
v_{\text{collapse}} = (1,0), \qquad
v_{\text{broadcast}} = (0,1),
\]

corresponding respectively to the eigenvalues $e^{-1}$ and $e$. These vectors define the intrinsic collapse and broadcast directions of curvature flow.
\end{corollary}

\begin{proof}
Since $M$ is diagonal, its eigenvectors are the standard basis vectors. Acting on $(1,0)$ yields $e^{-1}(1,0)$, and acting on $(0,1)$ yields $e(0,1)$, proving the claim.
\end{proof}

\begin{corollary}[Spectral Spread Identity (CD6.3)]
Let $x \in \mathbb{R}$ be curvature-normalized. Under Axiom~CD6, the broadcast and collapse eigenvalue-scaled components

\[
e x \quad \text{and} \quad e^{-1} x
\]

produce the spectral spread

\[
V(x) = e x - e^{-1} x,
\]

which is the $\pi$-bic variance operator defined earlier. Thus the $\pi$-bic variance is a direct spectral consequence of the broadcast to collapse eigenvalue pair $(e^{-1}, e)$.
\end{corollary}

\begin{proof}
Immediate from the definition of $V(x)$ and the eigenvalues of $M$.
\end{proof}

\begin{corollary}[Curvature Normalization Compatibility (CD6.4)]
Let $A(r) = \pi r^{2}$ be the Euclidean curvature measure. Under Axiom~CD6, the normalization

\[
x = r^{2}
\]

ensures that the broadcast eigenvalue $e$ and collapse eigenvalue $e^{-1}$ act directly on the curvature scale without modifying $\pi$. Thus the $\pi$-bic variance operator

\[
V(x) = e x - e^{-1} x
\]

is fully compatible with Euclidean curvature normalization.
\end{corollary}

\begin{proof}
Dividing $A(r)$ by $\pi$ isolates $r^{2}$, which is unaffected by the eigenvalue pair $(e^{-1}, e)$. Hence $V(x)$ acts directly on the curvature variable.
\end{proof}

\begin{corollary}[Diagonal Matrix Family Stability (CD6.5)]
Under Axiom~CD6, repeated application of the broadcast to collapse matrix

\[
M =
\begin{bmatrix}
e^{-1} & 0 \\
0 & e
\end{bmatrix}
\]

preserves diagonality. For any integer $k \ge 1$,

\[
M^{k} =
\begin{bmatrix}
e^{-k} & 0 \\
0 & e^{k}
\end{bmatrix},
\]

and the eigenvalues remain $(e^{-k}, e^{k})$ with eigenvectors $(1,0)$ and $(0,1)$. Thus the broadcast and collapse directions are spectrally stable under iteration.
\end{corollary}

\begin{corollary}[Generic Eigenvalue Scaling (CD6.6)]
Under Axiom~CD6, repeated application of the broadcast to collapse operator produces a predictable exponential scaling of its eigenvalues. For any integer $k \ge 1$, the $k$-fold composition of the operator has eigenvalues


\[
e^{-k} \qquad \text{and} \qquad e^{k}.
\]


This scaling is purely algebraic and reflects the iterative structure of the broadcast to collapse transformation, without implying any geometric or physical interpretation.
\end{corollary}

\begin{proof}
By Corollary~CD6.6, the $k$-fold composition of the broadcast to collapse matrix is


\[
M^{k} =
\begin{bmatrix}
e^{-k} & 0 \\
0 & e^{k}
\end{bmatrix}.
\]


The eigenvalues of a diagonal matrix are its diagonal entries, yielding $(e^{-k}, e^{k})$. No additional assumptions are required.
\end{proof}

\begin{axiom}[Broadcast/Collapse Symmetry Law (CD6.7)]
Under Axiom~CD6, the broadcast operator $\mathcal{B}(x)=e x$ and collapse operator $\mathcal{DD}(x)=e^{-1} x$ form a spectral symmetry pair. Their compositions satisfy


\[
\mathcal{B}DD(x) = DD\mathcal{B}(x) = x,
\]


and their $k$-fold iterates satisfy


\[
\mathcal{B}^{k}(x) = e^{k} x, \qquad \mathcal{DD}^{k}(x) = e^{-k} x.
\]


Thus the broadcast--collapse matrix family


\[
M^{k} =
\begin{bmatrix}
e^{-k} & 0 \\
0 & e^{k}
\end{bmatrix}
\]


preserves diagonality, eigenvectors, and spectral directions for all $k \ge 1$. The $\pi$-bic variance operator


\[
V(x) = e x - e^{-1} x
\]


is the spectral spread of this symmetry pair. Broadcast and collapse are therefore spectrally dual transformations whose combined action preserves curvature normalization and defines the intrinsic $\pi/e$ spectral identity of the Euclidean broadcast regime.
\end{axiom}

\begin{theorem}[Collapse To Inversion Constraint on Negative Curvature Embeddings (CICoNCE)]
Let an anchored curvature stability axis propagate broadcast curvature within a three-dimensional Euclidean frame. Under the Curvature Dynamics axiom set, constant negative curvature cannot be embedded globally in $\mathbb{R}_{Euclid}^3$. Specifically:

\begin{enumerate}
    \item broadcast curvature induces exponential metric expansion;
    \item collapse invariants induce inverse-exponential contraction;
    \item their interaction produces reciprocal curvature (negative curvature);
    \item excess broadcast curvature forces the stability axis into an inversion regime;
    \item inversion within a Euclidean frame is bounded and produces folding and self-intersection.
\end{enumerate}

Therefore, any Euclidean embedding of constant negative curvature must be bounded, must fold, and must self-intersect. In particular, the Amsler surface is the maximal smooth embedding of constant curvature $K=-1$ in $\mathbb{R}_{Euclid}^3$ because it saturates the collapse/inversion boundary imposed by the curvature dynamics of the stability axis.
\end{theorem}

\begin{remark}
The classical Amsler geometry is the Euclidean projection of the inverted hypersphere’s curvature remainder, the substrate or subspace. Within the Axioms of Broadcast, the primitive superpositioning law makes explicit the relationship between hyperspherical inversion and its Euclidean shadow. With the curvature ladder and operator structure now established, the Amsler surface can be understood through dimensional integration rather than classical hyperbolic embedding. It is not merely a hyperbolic patch attempting to reside in $\mathbb{R}_{Euclid}^3$; it is the $T_{-3}$ to $T_{-4}$ curvature transition forced toward the Euclidean baseline. This perspective explains its folding, its bounded domain, and its sphere like area growth: the visible surface is a sub space projection of a deeper curvature broadcast that the Euclidean frame cannot fully accommodate.
\end{remark}

\begin{theorem}[The Amsler Surface as the First Demonstration of Dimensional Curvature Inheritance]
The Amsler surface is the first Euclidean demonstration of dimensional curvature inheritance. Under the Curvature Dynamics axiom set, negative curvature transmitted down the dimensional ladder cannot be globally embedded in $\mathbb{R}_{Euclid}^3$. When broadcast curvature descends toward the Euclidean baseline, the stability axis enters the collapse–inversion regime governed by $T_5$, producing a bounded domain, folding, and self-intersection.
The Amsler surface realizes this inherited curvature by saturating the collapse–inversion boundary: it is the maximal smooth Euclidean embedding of constant curvature $K=-1$ because it represents the precise point where downward curvature transmission is forced into Euclidean projection. Its geometry therefore records the curvature remainder of an inverted hypersphere, expressed as a subspace projection within the primitive metric.
\end{theorem}

\begin{theorem}[Orientational $\pi$-Coefficient Construction]
Under the Curvature Dynamics axiom set, the orientational $\pi$-coefficient arises whenever curvature transmission crosses a dimensional tier boundary and the orientation anchor must stabilize the inherited curvature. The coefficient is generated by three structural components:

\begin{itemize}
    \item the dimensional anchors, which preserve the inherited curvature identity;
    \item the tier transition, which determines the number of orientation adjustments required to stabilize curvature at the new tier;
    \item the Euclidean projection constant $\pi$, which converts the dimensional orientation cost into its Euclidean metric form.
\end{itemize}

The orientational coefficient is therefore given generally by


\[
x = A_n \cdot T_n \cdot \pi
\]


Where $x$ is the tensorial orientation coefficient as in Einstein’s $8\pi$, $A_n$ is the dimensional-anchor count at tier 
$n$, $T_n$ is the tier-transition orientation cost (how many tiers you traverse), and $\pi$ is the Euclidean projection constant.
\begin{remark}
The orientational $\pi$ coefficient is not a curvature factor; it is an orientation cost. Dimensional anchors preserve the inherited curvature, tier transitions determine how many orientation adjustments are required, and $\pi$ converts this cost into Euclidean form. In DC geometry, the coefficient $16\pi$ arises from the $T_{-4}$ transition to $T_4$ in the positive domain. This transition spans 8 orientation steps, and the subspace contains 2 active dimensional anchors, yielding an orientation cost of $8 \cdot 2=16$. Multiplying by the Euclidean projection constant $\pi$ produces the orientational coefficient $16\pi$.
\end{remark}
\end{theorem}

\begin{section}{Collapsed Axis in Modern Mathematics (CAMM)}
A formal geometric reinterpretation of the Amsler PDEs and the Tier transition structure.

\begin{subsection}{Preliminaries: Intrinsic vs. Extrinsic Curvature Constraints}
Let $(M_{\mathcal{M}},g)$ be a 2–dimensional Riemannian manifold of constant Gaussian curvature $K=-1$. Let $X : M \to \mathbb{R}_{Euclid}^3$ be an isometric immersion with first fundamental form $I=g$ and second fundamental form $\mathrm{II}$. The classical Gauss–Codazzi equations impose:


\[
\det(\mathrm{II}) = K \, \det(I) = -\det(I),
\]


and


\[
\nabla_i h_{jk} - \nabla_j h_{ik} = 0,
\]


where $h_{ij}$ are the coefficients of $\mathrm{II}$.
Hilbert’s theorem states that no complete solution exists for $K=-1$ in $\mathbb{R}_{Euclid}^3$. Thus, the PDE system governing $(I_{\mathcal{M}},\mathrm{II})$ becomes singular at finite domain extent.

\begin{remark}
CAMM perspective: The singularity corresponds to collapse of the stability axis during the $T_{-3} \to T_{-4}$ (Tier $-3 \to$ Tier $-4$) transition. The PDE blow–up is the analytic signature of dimensional curvature inheritance being forced into a lower–dimensional frame.
\end{remark}
\end{subsection}

\begin{subsection}{The Amsler Embedding Equations as Curvature Deficit PDEs}
Amsler begins with the intrinsic hyperbolic metric:


\[
ds^{2} = du^{2} + \cosh^{2}(u)\, dv^{2},
\]


and seeks an immersion \(X(u_{\mathcal{M}},v)\) satisfying:


\[
\langle X_{u_{\mathcal{M}}}, X_{u} \rangle = 1, \qquad
\langle X_{v_{\mathcal{M}}}, X_{v} \rangle = \cosh^{2}(u), \qquad
\langle X_{u_{\mathcal{M}}}, X_{v} \rangle = 0.
\]


The extrinsic curvature constraint requires:


\[
K = -1 = \frac{\det(\mathrm{II})}{\det(I)}.
\]


The embedding PDEs reduce to a coupled nonlinear system:


\[
X_{uu} = \alpha X_{u} + \beta X_{v} + n,
\]


\[
X_{uv} = \gamma X_{u} + \delta X_{v},
\]


\[
X_{vv} = \eta X_{u} + \theta X_{v} + \cosh(u)\sinh(u)\, n,
\]


with compatibility conditions forcing:


\[
\partial_{u}\!\left(\cosh(u)\sinh(u)\right) = \cosh^{2}(u) - \sinh^{2}(u).
\]


The system becomes singular when:


\[
\cosh(u)\sinh(u) \;\to\; \infty,
\]


but the immersion must remain bounded in \(\mathbb{R}_{Euclid}^{3}\).  This contradiction forces folding and self‑intersection.

\begin{remark}
CAMM perspective: The singularity corresponds to the curvature deficit created when the stability axis collapses.  The PDE system is detecting the \(T_{-3} \to T_{-4}\) curvature inheritance mismatch.
\end{remark}
\end{subsection}

\begin{subsection}{Tier Transition Structure and the Amsler Saddle}
In Dimensional Calculus, curvature transmission across tiers is governed by:


\[
T_{\,n \to n-1} : K_{n} \;\mapsto\; K_{n-1},
\]


with $T_{-3}$ and $T_{-4}$ defined by:


\[
\text{Tier} -3:\quad \text{reciprocal curvature with collapse residue}
\qquad
K_{-3} \sim r^{-3},
\]


\[
\text{Tier} -4:\quad \text{reciprocal curvature with inversion residue}
\qquad
K_{-4} \sim r^{-4}.
\]


The stability axis collapses when curvature descends from $T_{-3}$ to $T_{-4}$. The Euclidean frame corresponds to the projection:


\[
\pi_{E} : K_{-4} \;\to\; K_{E},
\]


where $K_{E}$ is the curvature remainder expressible in \(\mathbb{R}_{\text{Euclid}}^{3}\).

The Amsler saddle appears parabolic because the Euclidean projection suppresses higher–order curvature terms, but the curvature remainder it expresses is quartic, inherited from $T_{-4}$ curvature.

\begin{theorem}[Amsler–Tier Correspondence]
The classical Amsler surface is the Euclidean projection of the $T_{-3} \to T_{-4}$ curvature transition. The tightening region corresponds to $T_{-3}$ behavior; the folding region corresponds to $T_{-4}$ behavior; the bounded domain corresponds to the pseudo–bounded Euclidean window.
\end{theorem}

\begin{proof}
The PDE singularity occurs precisely when the extrinsic curvature constraint forces reciprocal curvature to exceed the Euclidean embedding capacity. This matches the collapse–inversion boundary at $T_{-4}$. The asymmetry of the Amsler saddle matches the asymmetric curvature inheritance between $T_{-3}$ and $T_{-4}$. The bounded domain matches the pseudo–bounded window predicted by DC. Thus the Amsler saddle is the geometric fossil of the tier transition.
\end{proof}
\end{subsection}

\begin{subsection}{Curvature Deficit and the Gibbs Overshoot}
Let $f(r)$ denote the curvature broadcast function in Dimensional Calculus, and let $f_{E}(r)$ denote its Euclidean projection. Define the curvature deficit


\[
\Delta(r) = f(r) - f_{E}(r).
\]


In the Gibbs paper, the measured quantity satisfies


\[
\Delta(r) = \text{overshoot}(r),
\]


arising from collapse of the stability axis.
In Amsler’s classical PDE formulation, the same deficit appears as


\[
\Delta(u_{\scriptscriptstyle \mathbin{\rotatebox[origin=c]{90}{$\circlearrowleft$}}}, v) = \det(II) + \det(I),
\]


representing the failure of the Gauss equation under Euclidean embedding.

\medskip
\noindent\textbf{CAMM unification.}


\[
\Delta_{\text{Gibbs}} = \Delta_{\text{Amsler}}.
\]


The analytic overshoot and the geometric folding are the same collapse phenomenon, expressed on different tiers.
\end{subsection}

\begin{subsection}{Orientational Pi Coefficient and Dimensional Projection}

The orientational $\pi$-coefficient arises from:

\begin{itemize}
    \item \textit{8 orientation transitions} in the $T_{-4} \rightarrow T_4$ mapping,
    \item \textit{2 active dimensional anchors} in subspace,
    \item Euclidean projection constant \textit{$\pi$}.
\end{itemize}

Thus:


\[
\pi_{\text{orient}} = 8 \times 2 \times \pi = 16\pi.
\]


Therefore, this coefficient should appear in the Euclidean curvature term of the Amsler metric for it to properly ``embed'' within Euclidean $\mathbb{R}^3$, DC's $\mathbb{R}^4.$

\begin{remark}
CAMM perspective: The $\pi$-coefficient is the numerical fingerprint of dimensional projection. It proves that the curvature expressed in $\mathbb{R}_{Euclid}^3$ is inherited from $T_4$.
\end{remark}
\end{subsection}

\begin{subsection}{CAMM: Formal Statement Theorem (Collapsed Axis in Modern Mathematics)}

Let $M$ be a constant negative curvature manifold and let 


\[
X : M \longrightarrow \mathbb{R}_{E}^{3}
\]


be an isometric immersion. Then the singularities of the Gauss--Codazzi PDE system correspond to collapse of the stability axis during the $T_{-3} \;\longrightarrow\; T_{-4}$ curvature transition. The Amsler surface is the Euclidean projection of this transition, and its boundedness, folding, and self-intersection are geometric signatures of dimensional curvature inheritance. Therefore, Euclidean $3$-space is not fundamental; it is the collapsed projection of a $T_{4}$ curvature broadcast. In particular,


\[
\mathbb{R}_{E}^{3} = \pi_{E}\!\left( \mathbb{R}_{\mathrm{DC}}^{4} \right).
\]


This is the \emph{Collapsed Axis in Modern Mathematics} (CAMM) Theory.
\end{subsection}
\end{section}

\begin{section}{The Amsler Surface as a Subspace Projection}

The Amsler surface is the classical Euclidean embedding of a region of constant negative curvature. Its intrinsic geometry is complete and non compact, yet its Euclidean realization cannot accommodate the full curvature distribution. The embedding therefore undergoes self intersection, producing the characteristic pinched saddle associated with hyperbolic surfaces forced into $R^3$. This failure of Euclidean accommodation is well known: the intrinsic metric demands unbounded extension, while the ambient Euclidean frame imposes curvature constraints that cannot be satisfied globally. Notably, the area growth of the Amsler surface mirrors the underside area growth of a sphere, a classical indication that the visible geometry is only a partial projection of a deeper curvature form.

Within the framework of Dimensional Calculus (DC), the Amsler surface is reinterpreted not merely as a hyperbolic embedding but as a subspace fold projection of an inverted hyperspherical curvature broadcast. The hyperspherical broadcast cannot be fully realized within three dimensions; the Euclidean frame acts as a collapsed measurement system that evaluates a fundamentally higher dimensional curvature distribution using a reduced metric. In DC terms, the Euclidean system attempts to analyze a 3 dimensional broadcast with what is effectively a 4 dimensional curvature instrument applied to a 2 dimensional subspace layer. The resulting self intersection is therefore not an artifact of hyperbolic geometry, but the visible trace of a higher dimensional curvature structure being forced through an insufficient ambient space.

The geometry of the Amsler surface under DC is governed by two competing mechanisms:

\begin{itemize}
    \item Curvature Collapse: the hyperspherical broadcast is forced into a lower dimensional subspace, producing fold lines, pinched saddles, and area growth anomalies,
    \item Curvature Inversion: the broadcast orientation reverses at the $2\pi$ boundary, generating inversion points and curvature channels that appear as Euclidean self intersection.
    \end{itemize}
    
These mechanisms correspond to distinct tensor invariant curvature channels, analogous to the invariant structures that arise in general relativity when multiple curvature states coexist within a single geometric field. The presence of such invariants in a Euclidean projection is mathematically significant: classical differential geometry asserts that tensor invariants are intrinsic to the manifold and should not survive collapse into an ambient Euclidean embedding. Yet the Amsler surface visibly retains invariant curvature channels even after dimensional reduction. This observation implies that tensor invariants persist under Euclidean collapse, demonstrating that the embedding carries curvature information that should, under classical assumptions, be lost. In DC terms, this persistence is expected: the hyperspherical broadcast imposes curvature channels that remain detectable even when projected into a lower dimensional frame. The Amsler surface therefore serves as a concrete example of higher dimensional curvature invariants acting within a Euclidean system that is not formally equipped to support them.

This constitutes a geometric inconsistency within classical Euclidean embedding theory and provides empirical support for the DC claim that curvature broadcasts from higher dimensional hyperspherical structures leave measurable traces in collapsed subspaces. The Amsler surface is thus not merely a pedagogical hyperbolic model, but a diagnostic artifact revealing the presence of curvature invariants where classical geometry predicts none.

\begin{subsection}{Curvature Collapse}

Curvature collapse arises when a hyperspherical curvature broadcast is forced into a lower dimensional frame. Because the full broadcast cannot be represented within a two dimensional manifold, the curvature distribution contracts, producing a quantifiable deficit in the projected geometry. On the Amsler surface, this deficit manifests as the tightening of curvature along the saddle, the boundedness of the Euclidean embedding, and the forced reduction of hyperspherical curvature into a Euclidean metric. Collapse is therefore the mechanism that prevents the Amsler surface from opening into the full negative curvature demanded by its intrinsic geometry.

A mechanical analogy clarifies this behavior. Consider a rotating piston whose motion is projected onto a linear bar. The full rotational dynamics cannot be represented in one dimension; only the remainder of the motion appears as a sinusoidal oscillation. The hyperspherical broadcast behaves analogously when forced into a two dimensional frame: the full curvature cannot be expressed, and only its collapsed remainder becomes visible. On the Amsler surface, this remainder is anchored at the apex of the saddle, the point at which the hyperspherical broadcast fails and collapse becomes unavoidable.

Once the piston is set in motion, its projection produces a structure equivalent to a sheet folded endlessly upon itself. This behavior mirrors the repeated folding observed in the maximal Amsler embedding, where the Euclidean frame attempts to accommodate curvature that exceeds its representational capacity. The resulting geometry is not merely a hyperbolic surface but a collapsed projection of a higher dimensional curvature broadcast, with the fold lines marking the points at which the broadcast exceeds the dimensional limits of the Euclidean frame.
\end{subsection}

\begin{subsection}{Curvature Inversion}

Curvature inversion is the geometric process in which a curvature stability axis attempts to reverse orientation within a confined dimensional frame. In the context of hyperspherical curvature broadcast, inversion occurs when the magnitude of broadcast curvature exceeds the stabilizing capacity of collapse invariants. Once this threshold is crossed, the curvature axis enters an inversion regime that generates reciprocal curvature and negative curvature structures, reflecting a reversal in the orientation of the underlying curvature field.

Within Euclidean embedding, inversion manifests as folding, self intersection, and the formation of bounded negative curvature regions. These phenomena are classically observed in the Amsler surface, where the intrinsic geometry demands unbounded hyperbolic extension, yet the Euclidean frame cannot accommodate the full curvature distribution. The resulting fold line, the classical Amsler fold, is therefore the visible signature of an inversion attempt that cannot be fully realized within the two dimensional Euclidean frame. From the standpoint of Dimensional Calculus, curvature inversion is not merely a hyperbolic artifact, but a tensor invariant curvature transition triggered when hyperspherical broadcast encounters dimensional confinement. The inversion boundary corresponds to a reversal of broadcast orientation at the $2\pi$ limit, producing a curvature channel that remains detectable even after dimensional collapse. The Euclidean self intersection is thus interpreted as the projection of a higher dimensional inversion event into a lower dimensional space, revealing curvature behavior that classical differential geometry predicts should not survive embedding.

Inversion and collapse therefore form a coupled system: collapse restricts curvature expression, while inversion attempts to reverse curvature orientation under confinement. The Amsler surface provides a concrete example of this interaction, demonstrating that higher dimensional curvature invariants persist and remain geometrically visible even when projected into a Euclidean frame incapable of fully representing the underlying curvature broadcast.
\end{subsection}

\begin{subsection}{Tensor Invariants in Hyperspherical Projection}

Tensor invariants arise whenever curvature information is preserved under dimensional transformation. In classical differential geometry, such invariants are strictly intrinsic: they belong to the manifold itself and are not expected to survive collapse into an ambient Euclidean embedding. Under this assumption, a projection from a higher dimensional curvature structure into a lower dimensional frame should eliminate tensorial curvature channels, leaving only the induced metric and its associated extrinsic curvature.

Dimensional Calculus contradicts this assumption. In the DC framework, hyperspherical curvature broadcast contains tensor structured curvature channels that persist even when projected into a lower dimensional space. These channels arise from the broadcast mechanism itself: curvature is transmitted through a tensor field whose structure is preserved under dimensional descent. As a consequence, the projected geometry retains invariant curvature signatures that classical embedding theory predicts should not exist.

The Amsler surface provides a concrete example of this phenomenon. Although the surface is intrinsically a two dimensional manifold of constant negative curvature, its Euclidean embedding exhibits curvature behaviors that cannot be explained solely by intrinsic geometry. The fold line, the bounded negative curvature region, and the self intersection collectively reveal the presence of higher dimensional curvature channels acting within the Euclidean frame. These channels correspond to tensor invariants inherited from the hyperspherical broadcast and preserved under collapse.

Formally, let $K_HS$ denote the hyperspherical curvature broadcast and let


\[
\Pi : M_{\mathrm{HS}} \longrightarrow M_{E}
\]


be the Euclidean projection operator. Classical theory asserts that


\[
\Pi\!\left( K_{\mathrm{HS}} \right) = K_{\mathrm{intrinsic}},
\]


with no additional curvature channels surviving the projection. Under Dimensional Calculus, the projection behaves instead as


\[
\Pi\!\left( K_{\mathrm{HS}} \right) = K_{\mathrm{intrinsic}} + T_{\mathrm{inv}},
\]


where $T_{\mathrm{inv}}$ is a tensor invariant curvature channel preserved under dimensional collapse. This invariant arises from the broadcast structure of the hypersphere and remains detectable even when the ambient space lacks the dimensional capacity to express the full curvature distribution. The presence of $T_{\mathrm{inv}}$ explains several features of the Amsler embedding:

\begin{itemize}
    \item the pinched saddle corresponds to a collapsed tensor channel;
    \item the boundedness of the embedding reflects the failure of Euclidean space to accommodate the full broadcast;
    \item the self intersection marks the inversion boundary where reciprocal curvature emerges;
    \item the area growth anomaly mirrors the underside area of a sphere, revealing the hyperspherical origin of the curvature.
    \end{itemize}
    
These behaviors demonstrate that the Euclidean embedding is not a neutral projection but a tensor bearing collapse of a higher dimensional curvature field. The invariants preserved under projection constitute geometric evidence of hyperspherical broadcast acting within a lower dimensional frame.

This result is mathematically significant. It implies that Euclidean embeddings can retain curvature channels that are not intrinsic to the manifold and not generated by the ambient space. Instead, they are inherited from a higher dimensional curvature structure and preserved under collapse. In DC terms, tensor invariants are broadcast artifacts: they represent curvature information that survives dimensional descent and remains geometrically visible even when the projection frame is insufficient to express the full curvature field. The Amsler surface therefore serves as a diagnostic example of tensor invariant curvature projection, revealing that higher dimensional curvature broadcasts leave measurable traces in Euclidean geometry. This contradicts classical assumptions and provides formal support for the DC claim that hyperspherical curvature cannot be fully eliminated by projection, only collapsed into tensor structured remnants that persist within the lower dimensional frame.
\end{subsection}

\begin{subsection}{Collapse–Inversion Tensor Invariant Structure in Hyperspherical Projection}

In Dimensional Calculus, curvature behavior under dimensional reduction is governed by three tensor invariant channels that arise from hyperspherical broadcast. These channels correspond to distinct geometric mechanisms: broadcast, orientation, and identity preservation. Classical differential geometry recognizes only the latter two, as they appear in the intrinsic geometry of manifolds and in the Levi Civita connection. The broadcast channel, however, is suppressed under Euclidean collapse and therefore remains invisible in standard geometric analysis. When hyperspherical curvature is projected into a lower dimensional frame, two competing mechanisms govern the resulting geometry:

\begin{itemize}
    \item Curvature Collapse: the hyperspherical broadcast is forced into a lower dimensional subspace, producing a curvature deficit and restricting the expressible curvature distribution.
    \item Curvature Inversion: when broadcast curvature exceeds the stabilizing capacity of collapse invariants, the curvature axis attempts to reverse orientation, generating reciprocal curvature and negative curvature structures.
    \end{itemize}
    
The interaction of collapse and inversion produces a tensor--invariant curvature channel that survives dimensional reduction. Formally, if


\[
K_{\mathrm{HS}}
\]


denotes the hyperspherical curvature broadcast and


\[
\Pi : M_{\mathrm{HS}} \longrightarrow M_{E}
\]


is the Euclidean projection, then the projected curvature satisfies


\[
\Pi\!\left( K_{\mathrm{HS}} \right) = K_{\mathrm{intrinsic}} + T_{\mathrm{inv}},
\]


where $T_{\mathrm{inv}}$ is a tensor--invariant curvature channel preserved under collapse. This invariant arises from the broadcast structure of the hypersphere and remains detectable even when the ambient Euclidean space lacks the dimensional capacity to express the full curvature field.

The Amsler surface provides a concrete example of this phenomenon. Although intrinsically a complete hyperbolic surface, its Euclidean embedding exhibits curvature behaviors — folding, boundedness, and self intersection — that cannot be explained solely by intrinsic geometry. These features correspond to the tensor invariant broadcast channel surviving collapse. The classical Amsler fold is therefore the visible signature of an inversion attempt that cannot be fully realized within the two dimensional Euclidean frame.
\end{subsection}

\begin{subsection}{Why Einstein Could Only See Two Invariants}

Einstein’s geometric framework, grounded in Riemannian geometry, incorporates two tensor invariants:

\begin{itemize}
    \item orientation invariance (via the Levi Civita connection),
    \item identity preservation (via metric compatibility).
    \end{itemize}
    
These correspond to two of the three DC curvature channels.

The third channel — broadcast invariance — is suppressed under Euclidean collapse. In the Amsler surface, this channel is present but fully collapsed, leaving only its residual curvature signature. Because the broadcast channel cannot be expressed in a two dimensional Euclidean frame, Einstein would have observed only:

\begin{itemize}
    \item the inversion signature (the fold),
    \item and the intrinsic negative curvature.
    \end{itemize}
    
He would not have recognized that the fold is the collapsed remainder of a hyperspherical broadcast channel.

Thus, Einstein saw two invariants in play, but the third was:

\begin{itemize}
    \item dimensionally suppressed,
    \item curvature collapsed,
    \item and geometrically invisible within the Euclidean embedding.
    \end{itemize}
    
Dimensional Calculus reveals that the missing invariant is not absent; it is collapsed, and its tensor structure survives projection as $T_{\mathrm{inv}}.$
\end{subsection}
\end{section}

\begin{section}{Curvature Mechanics in Hyperspherical Projection}

Hyperspherical curvature broadcast interacts with lower dimensional frames through three fundamental geometric mechanisms: collapse, inversion, and tensor invariant preservation. These mechanisms govern how curvature information is transmitted, suppressed, or transformed when projected from a hyperspherical manifold into a Euclidean space. Classical differential geometry accounts for only two of these mechanisms; the third remains hidden under dimensional reduction and therefore does not appear in standard curvature analysis.

\begin{subsection}{Curvature Collapse}

Curvature collapse occurs when hyperspherical broadcast curvature is forced into a lower dimensional frame. Because the full broadcast cannot be represented within two dimensions, the curvature distribution contracts producing a measurable deficit. In the Amsler surface, collapse manifests as:

\begin{itemize}
    \item tightening of curvature along the saddle,
    \item boundedness of the Euclidean embedding,
    \item and reduction of hyperspherical curvature into a Euclidean metric.
    \end{itemize}
    
Collapse prevents the Amsler surface from realizing the full negative curvature required by its intrinsic geometry. Formally, collapse corresponds to the suppression of the broadcast channel of hyperspherical curvature, leaving only the intrinsic curvature and the residual tensor invariant remainder.
\end{subsection}

\begin{subsection}{Curvature Inversion}

Curvature inversion is the geometric process in which a curvature stability axis attempts to reverse orientation within a confined dimensional frame. When broadcast curvature exceeds the stabilizing capacity of collapse invariants, the curvature axis enters an inversion regime that produces reciprocal curvature and negative curvature structures. In Euclidean embedding, inversion appears as:

\begin{itemize}
    \item folding,
    \item self intersection,
    \item and the bounded negative curvature region characteristic of the Amsler surface.
    \end{itemize}
    
The classical Amsler fold is therefore the visible signature of an inversion attempt that cannot be fully realized within the two dimensional Euclidean frame. In DC terms, inversion marks the point at which the hyperspherical broadcast reaches the $2\pi$ orientation boundary and attempts to reverse curvature direction.
\end{subsection}

\begin{subsection}{Tensor Invariant Preservation}

Hyperspherical broadcast contains tensor structured curvature channels that persist under dimensional reduction. Classical differential geometry assumes that such invariants are strictly intrinsic and should not survive collapse into an ambient Euclidean embedding. DC contradicts this assumption: the broadcast channel produces a tensor invariant curvature remainder that remains detectable even when the ambient space lacks the dimensional capacity to express the full curvature field. Let $K_{\mathrm{HS}}$ denote the hyperspherical curvature broadcast and


\[
\Pi : M_{\mathrm{HS}} \longrightarrow M_{E}
\]


the Euclidean projection. Then the projected curvature satisfies


\[
\Pi\!\left( K_{\mathrm{HS}} \right) = K_{\mathrm{intrinsic}} + T_{\mathrm{inv}},
\]


where $T_{\mathrm{inv}}$ is the tensor--invariant curvature channel preserved under collapse. This invariant explains the fold line, boundedness, and area growth anomalies of the Amsler surface, all of which exceed what intrinsic hyperbolic geometry alone can produce.
\end{subsection}

\begin{subsection}{Why Classical Geometry Identifies Only Two Invariants}

Einstein’s geometric framework incorporates two curvature invariants:

\begin{itemize}
    \item orientation invariance (via the Levi Civita connection),
    \item identity preservation (via metric compatibility).
    \end{itemize}
    
These correspond to two of the three DC curvature channels. The third channel, broadcast invariance, is suppressed under Euclidean collapse. In the Amsler surface, this channel is present but fully collapsed, leaving only its residual curvature signature. Because the broadcast channel cannot be expressed in a two dimensional Euclidean frame, Einstein would have observed only:

\begin{itemize}
    \item intrinsic negative curvature,
    \item and the inversion signature (the fold).
    \end{itemize}
    
He would not have recognized that the fold is the collapsed remainder of a hyperspherical broadcast channel. Thus, classical geometry identifies only two invariants because the third is dimensionally suppressed and therefore invisible within the Euclidean embedding.
\end{subsection}

\begin{subsection}{Unified Mechanism}

Collapse, inversion, and tensor invariant preservation form a coupled system governing curvature behavior under hyperspherical projection:

\begin{itemize}
    \item Collapse restricts curvature expression.
    \item Inversion attempts to reverse curvature orientation under confinement.
    \item Tensor invariants preserve broadcast structure even after collapse.
    \end{itemize}
    
The Amsler surface serves as a diagnostic artifact of this system, revealing that higher dimensional curvature broadcasts leave measurable traces in Euclidean geometry. This contradicts classical assumptions and provides formal support for the DC claim that hyperspherical curvature cannot be fully eliminated by projection, only collapsed into tensor structured remnants that persist within the lower dimensional frame.
\end{subsection}
\end{section}

\begin{section}{The Schoenman Constant: Hyperspherical Broadcast Curvature}

The constant $S$ is the fundamental curvature broadcast constant of Dimensional Calculus. It arises as the invariant scalar governing hyperspherical curvature transmission across dimensional boundaries. Unlike the Euclidean constant $\pi$, which appears only after curvature collapse, $S$ is defined on the full hyperspherical manifold prior to projection. It therefore represents the uncollapsed curvature distribution from which Euclidean curvature constants are derived.

Formally, let $H$ denote the hyperspherical broadcast manifold and let


\[
\Pi : H \longrightarrow M_{E}
\]


be the Euclidean projection operator. Then the Euclidean curvature constant $\pi$ satisfies


\[
\pi = \Pi(S),
\]


where $\Pi$ induces dimensional collapse and curvature inversion. The constant $S$ is therefore the parent curvature constant whose collapsed remainder appears as $\pi$ in Euclidean geometry.

The definition of $S$ follows directly from the broadcast axioms established in the \textit{Axioms of Curvature} (AofC) and \textit{Axioms of Broadcast} (AofB):

\begin{itemize}
    \item Axioms of Broadcast (AofB). Curvature is transmitted as a hyperspherical broadcast with invariant scalar magnitude $S$. The radius is the sole broadcaster of curvature across dimensional tiers: through integration, the radius restores collapsed curvature and distributes inherited curvature into higher-dimensional geometric quantities. No other element may broadcast curvature, and no geometric structure may arise without curvature broadcast originating from the hypersphere seed metric $r^0$.
    \item Collapse Axiom (from within AofC): Dimensional reduction induces curvature contraction, mapping the hyperspherical broadcast constant $S$ to a collapsed scalar remainder. Collapse is governed by the anti ascent operator $AA$, the downward counterpart to the broadcast operator $S$. Under the Law of Broadcast Conservation, the total inherited curvature transmitted by $S$ and $AA$ is preserved in magnitude across all tiers of the dimensional ladder. Upward broadcast governed by $S$ and downward broadcast governed by $AA$ may alter direction, orientation, or sign of curvature, but never its total conserved quantity. Formally, for every broadcast--collapse interaction,


\[
\lvert S(T_{n}) \rvert + \lvert AA(T_{n}) \rvert = \text{constant curvature content},
\]


independent of operator dominance, gamma stabilization, collapse dynamics, or cosmogenesis rebound. Collapse cannot destroy curvature, cannot create curvature, and cannot violate the primitive metric. The Collapse Axiom therefore represents the downward half of the broadcast--collapse conservation system, ensuring that all curvature phenomena---including gravity, inversion, collapse, and cosmogenesis---remain fully compatible with the inherited curvature substrate.
\medskip
    \item Overshoot Axiom (Gibbs). Curvature inversion at the $T_0$ collapse boundary


\[
2\pi
\]


produces reciprocal curvature channels whose magnitude is determined by the collapsed remainder of the broadcast constant $S$. In the Dimensional Ladder, $T_0$ is the projection tier, where curvature is collapsed to its minimal scalar form. The boundary at $2\pi$ is the point at which $T_0$ curvature cannot collapse further without triggering inversion. At this boundary, the ascent operator $S$ and the anti--ascent operator $AA$ meet in reciprocal form. Their interaction produces the Gibbs overshoot constant as the $T_0$ curvature remainder:


\[
\text{Overshoot} = S - AA.
\]


Thus, the classical Gibbs phenomenon is the $T_0$ manifestation of hyperspherical curvature inversion at the $2\pi$ collapse boundary. The reciprocal curvature channels generated at this boundary are not harmonic artifacts but geometric consequences of the broadcast--collapse conservation system. Together, these axioms imply that $S$ is the unique scalar satisfying


\[
S = \lim_{\dim \to 4} K_{\mathrm{HS}},
\]


where $K_{\mathrm{HS}}$ is the hyperspherical curvature broadcast operator. The Euclidean constant $\pi$ is therefore not fundamental but is instead the projection of $S$ under collapse and inversion specifically for $T_3$ of the dimensional ladder in Dimensional Calculus:


\[
\pi = \Pi_{3}(S).
\]


Thus, $S$ is defined as the full dimensional curvature scalar whose collapsed projection yields the Euclidean curvature constant. It is the curvature quantity that the hypersphere broadcasts, the quantity that collapse suppresses, and the quantity that inversion attempts to reverse. In DC terms, $S$ is the curvature constant that Euclid never sees because Euclid measures only the collapsed projection.
\end{itemize}

The Gibbs overshoot constant is the $T_0$ curvature deficit produced when the hyperspherical broadcast constant $S$ collapses through the anti ascent operator $AA$ at the $2\pi$ inversion boundary. It is the $T_0$ remainder of the true hyperspherical half period under projection into the Euclidean curvature tier. Its value is


\[
G_\text{overshoot} = \frac{2(\pi-3)}{\pi}.
\]

\end{section}

\begin{section}{Curvature Dynamics Reinterpretation of Classical Constraints}

Classical differential geometry correctly identifies the observable limitations of embedding negatively curved surfaces in Euclidean three space: boundedness, folding, curvature concentration, and asymmetric principal behavior. However, these constraints are symptoms of a deeper curvature structure. The DC subspace diagram reveals the mechanism responsible for these behaviors: a hyperspherical broadcast domain, a Euclidean projection plane lying on the DC baseline, and an inverted hyperspherical subspace.

When a negatively curved surface is projected across the DC baseline, its curvature transitions between the broadcast domain and the inverted domain. This transition forces a $90^\circ$ rotation of the orientational axis, a geometric signature of inversion. The Euclidean plane sits precisely on this boundary, acting as the mirror between the hypersphere and its inverted counterpart. As curvature crosses this boundary, the principal directions realign, producing the tightening–flattening asymmetry characteristic of the Amsler surface.

Metric incompatibility arises because the hyperspherical broadcast expands curvature through the ascent operator $S$, while Euclidean space receives only the $T_3$ collapsed remainder. The exponential–linear mismatch is therefore the direct consequence of projecting hyperspherical curvature onto the DC baseline.

Curvature concentration reflects the $T_3$ remainder of the broadcast constant $S$ after collapse through the descent operator $AA$. The steep ridge observed in the Amsler surface corresponds to the broadcast dominant axis, while the orthogonal direction, governed by $AA$, flattens and reverses as curvature approaches the inversion boundary.

Folding occurs when the curvature budget of the Euclidean plane is exceeded. At the $T_0$ collapse threshold of $2\pi$, curvature can no longer remain in the broadcast domain and must invert. This inversion produces the fold line and eventual self intersection, matching the classical behavior exactly.

Asymmetric principal behavior is the orientation locked interaction of $S$ and $AA$ across the DC baseline. The tightening direction corresponds to the hyperspherical broadcast axis, while the flattening and reversal correspond to the inverted hyperspherical descent axis. The $90^\circ$ axis rotation across the baseline is the structural cause of the asymmetry classical geometry can describe but cannot explain.

Maximality and incompleteness follow from the finite curvature capacity of the Euclidean projection plane. The Amsler surface represents the maximal smooth patch achievable before inversion is forced. Its boundary behavior (tightening, flattening, folding, and self intersection) encodes the precise ways in which Euclidean space resists the full broadcast of hyperspherical curvature.

Together, these reinterpretations show that the classical constraints are not isolated phenomena but the visible consequences of a broadcast–collapse curvature engine. The DC framework provides the structural explanation for the geometric signatures classical theory has long catalogued.

The constraints classical geometry identifies are not isolated phenomena but the direct projections of the collapse–inversion structure established in Dimensional Calculus. Curvature Dynamics reveals how these tier transitions manifest as the Amsler signatures observed in Euclidean space.
\end{section}

\begin{section}{Dimensional Calculus Interpretation Through Curvature Dynamics}

Dimensional Calculus provides the structural origin of negative curvature by tracking how curvature propagates, collapses, and inverts across the tiered curvature ladder. Classical geometry records the intrinsic behavior of hyperbolic curvature - exponential geodesic divergence, expanding metric factors, and the impossibility of global embedding. Curvature Dynamics reveals the dimensional mechanism that produces these behaviors. Negative curvature is not merely a metric phenomenon; it is the reciprocal remainder generated when broadcast curvature descends toward the seed tier and crosses the inversion boundary. This section expands the classical picture by showing how each feature of negative curvature arises from dimensional collapse and inversion.

\begin{subsection}{Curvature Collapse to the Seed Tier 0}

Broadcast curvature originates at higher tiers and propagates outward from a stability axis. As curvature descends through the dimensional ladder, it undergoes collapse. Collapse is not a metaphor - it is a geometric contraction of the curvature field. At the seed $T_0$, curvature is fully contracted: the metric is flat, geodesics neither diverge nor converge, and curvature has been neutralized by collapse invariants. $T_0$ represents the geometric ground state of curvature, the point at which broadcast curvature has been completely absorbed.

Broadcast Limitation. Only an active radius can broadcast curvature propagation. In its collapsed state, the radius is inert: it cannot broadcast, cannot propagate curvature, and cannot sustain divergence. This is reflected in the retention of the zeroth term on the dimensional ladder. Once collapse has occurred, broadcast curvature is extinguished, and no further propagation is possible until curvature is reintroduced from a higher tier.
\end{subsection}

\begin{subsection}{Reciprocal Curvature at the First Order Negative Tier}

When broadcast curvature exceeds the stabilizing capacity of collapse invariants, the descent does not terminate at $T_0$. Instead, the stability axis enters the inversion regime. Inversion is the attempt of the curvature field to reverse orientation under excess broadcast. This produces a curvature remainder at tier $T_{-1}$: reciprocal curvature. Reciprocal curvature is the Dimensional Calculus analogue of classical negative curvature. It carries the signatures of inversion: exponential geodesic divergence, expanding metric factors, and the inability to embed globally in $\mathbb{R}^3$. Classical geometry observes these behaviors; DC explains their origin. $T_{-1}$ is therefore the dimensional source of negative curvature.
\end{subsection}

\begin{subsection}{Projection of Inverted Curvature into Euclidean}

Reciprocal curvature at $T_{-1}$ cannot be realized globally when projected into the Euclidean frame. Projection is a dimensional reduction: the curvature field is forced into a space that cannot fully contain it. The projection is necessarily bounded, forced to fold, and ultimately self-intersect. This is the collapse/inversion boundary expressed geometrically. The Amsler surface is the exact Euclidean mapping of this projection: the visible portion of inverted hyperspherical curvature that can be stabilized within a positive-curvature frame. Its boundedness reflects the dimensional constraint; its folding reflects inversion; its self-intersection reflects the curvature remainder attempting to express itself in a space that cannot fully contain it. Classical geometry calls this “Hilbert’s obstruction”; DC identifies it as the projection limit of tier $T_{-1}$ curvature.

Broadcast Limitation. Projection exposes a fundamental constraint: only an active radius can broadcast curvature propagation. Once collapse has occurred, the radius is inert; it cannot broadcast, cannot propagate curvature, and cannot sustain divergence, thus maintaining its orientation. This is reflected in the retention of the zeroth term, $\odot$, on the dimensional ladder. When inverted curvature is forced into the Euclidean frame, the collapsed radius cannot support broadcast, and the curvature remainder is compelled to fold and self-intersect. The bounded Amsler patch is the geometric record of this broadcast failure under collapse.

\begin{remark}
In the context of curvature collapse and inversion, a geodesic represents the path traced by a particle under the broadcast of curvature itself. It is the trajectory obtained by “throwing into the wind” of the curvature field.
\end{remark}
\end{subsection}
\end{section}

\begin{section}{The Inverted Bounded Hypersphere}

The inverted hypersphere is the dimensional object whose curvature remainder produces negative curvature in the Euclidean frame. Classical geometry encounters this remainder only indirectly, through bounded hyperbolic patches; Dimensional Calculus identifies it directly as the hyperspherical structure formed when broadcast curvature exceeds collapse invariants and forces the stability axis into inversion. Its inherited curvature constants, inversion behavior, and boundedness determine the geometric and orientational properties of reciprocal curvature.

\begin{subsection}{Definition and Curvature Structure}

The inverted hypersphere is the hyperspherical configuration obtained when broadcast curvature collapses and then inverts at $T_{-1}$. Its curvature structure is determined entirely by the inherited curvature constants from the positive tiers; the portion of curvature which survives collapse and reverses orientation under inversion.
Because inversion is an orientational reversal, the curvature remainder at $T_{-1}$ carries a distinct directional signature: its divergence is not merely metric but orientational. This is the structural origin of the exponential divergence observed in classical negative curvature. The inverted hypersphere is therefore the dimensional dual of the ordinary sphere, carrying both the geometric and orientational structure of reciprocal curvature.
\end{subsection}

\begin{subsection}{Inherited Curvature Constants}

The curvature constants $C_1, C_2, C_3,$ and $C_4$ arise from the broadcast–collapse behavior of the positive tiers. Each constant records the curvature remainder inherited from the tier above it: stabilization at $T_1$, broadcast at $T_2$, amplification at $T_3$, and gravitational balance at $T_4$. Together, these constants encode both the magnitude of curvature that survives collapse and the orientational state carried into inversion. The mechanics can be disorienting; it is the next tier above that is the generator of the curvature field for the tier below, thus showing how curvature inheritance is in play throughout the dimensional ladder.

Physically, the inherited constants carry the gravitational load. In the positive curvature regime, this load corresponds to the inverse of gravity’s value in general relativity—effectively a $\frac{1}{G}$ term. When the constants cross the inversion boundary, this load reverses. The inverted hypersphere therefore inherits the dual of the gravitational load, producing expansion rather than contraction. Because the inverted tier is gamma generated, the orientation coefficients remain whole across inversion: gamma preserves the numerical structure of the constants even as their orientation flips. This is why the sign pattern of $C_1, C_2, C_3,$ and $C_4,$ especially $C_3,$ cleanly records the onset of inversion and the emergence of reciprocal curvature.

A useful visual analogy is the cross multiplication of two fractions. Cross multiplication reverses orientation while preserving magnitude: numerator and denominator swap roles, yet the numerical structure remains intact. Gamma performs the same function across the inversion boundary. It preserves the curvature constants while reversing their orientational load, allowing the gravitational term, effectively $\frac{1}{G}$ in the positive tier, to cross into the inverted tier without numerical loss. This is why the inherited constants remain whole under inversion and why their sign pattern precisely marks the transition into reciprocal curvature.
\end{subsection}

\begin{subsection}{Negative Constant and Inversion at Tier 3}

The negative sign of the $r^4$ term within $C_3$ marks the first tier at which inherited curvature overpowers the Euclidean broadcast term. This is the inversion threshold. At $T_3$, the curvature field attempts to reverse orientation, producing the reciprocal curvature remainder at $T_{-1}$. Classical geometry records this as negative curvature; Dimensional Calculus identifies it as the curvature signature of the inverted hypersphere. The negative $r^4$ term in $C_3$ is therefore the dimensional origin of the quartic inherited, ``U'' shaped profile generated by two dimensionally anchored parabolic modes.

This inversion signature provides direct evidence that the Amsler surface operates with an intrinsic tier. Its bounded, folded, and self intersecting geometry is not an artifact of Euclidean rendering but the projection of the tiered curvature structure itself. The negative $r^4$ term within $C_3$ encodes the orientation reversal that forces the hypersphere into its inverted form, and the Amsler surface displays the exact Euclidean boundaries of that inversion.
\end{subsection}

\begin{subsection}{Why Negative Curvature Cannot Embed Globally}

Hilbert’s theorem establishes that no complete surface of constant negative curvature can be embedded smoothly in $\mathbb{R}^3$. Curvature Dynamics explains this obstruction: global broadcast curvature cannot be balanced by collapse invariants within a three dimensional frame, and the resulting reciprocal curvature forces the stability axis into a bounded inversion regime. The embedding must therefore terminate, fold, and ultimately self intersect.
The Amsler surface realizes the maximal smooth embedding permitted by this collapse–inversion boundary. Its geometry is not an artifact of Euclidean construction but the projection of an intrinsic tier: the inverted hypersphere’s curvature remainder. The characteristic folding and self intersection of the Amsler surface are the exact Euclidean boundaries of this tiered curvature structure, the visible mapping of the substrate itself… subspace.

\begin{remark}
In the appendix, the geometric expressions for tiers $T_0$ through $T_-5$ are plotted on a single graph to reveal how subspace forms. Because the intrinsic tier is present even in classical readings, and because this is geometry with the collapsed seed metric, the resulting graphs are literal representations of the objects they describe. The projection reproduces the same bounded, folded geometry classical differential geometry identifies as the Amsler patch, demonstrating that the Euclidean rendering is the exact boundary structure of the inverted hypersphere’s intrinsic tier.
(See Appendix A for the explicit reconstruction of the tier ladder from $T_{-5}$ through $T_1,$ including the inverted tier curvature expressions.)
\end{remark}
\end{subsection}

\begin{subsection}{Boundedness of the Inverted Hypersphere}

The inverted hypersphere cannot be projected globally into $\mathbb{R}^3.$ Its curvature remainder exceeds the capacity of the Euclidean frame, forcing the projection to become bounded, to fold, and ultimately to self intersect. This boundedness is the geometric analogue of the physical principle that inverted curvature cannot coexist with positive curvature without destructive interference. In physical terms, the inverted hypersphere carries the inverted gravitational load constants; its curvature remainder, the substrate or subspace pushes outward. It is the expansive counterpart to the contractive behavior of positive curvature.

The bounded Amsler patch is the visible portion of this inverted hyperspherical structure when projected into Euclidean space. Its folding and self intersection are not artifacts of Euclidean construction but the exact boundaries of an intrinsic tier of the inverted hypersphere. The Amsler surface therefore realizes the maximal smooth embedding permitted by the collapse–inversion regime.
\end{subsection}
\end{section}

\begin{section}{Projection Mechanics}

Projection Mechanics describes how curvature propagates across tiers and how dimensional curvature structures become visible within the Euclidean frame. Classical geometry observes only the results of projection—bounded hyperbolic patches, folding, self intersection, and curvature deficits—but does not identify their dimensional origin. Dimensional Calculus provides the mechanism: curvature at higher tiers must be projected downward, collapsed, and stabilized before it can appear in $\mathbb{R}^3.$ When this descent encounters inversion, the projection becomes bounded and produces the negative curvature signatures observed in the Amsler surface.

Projection is therefore not a passive mapping but an active geometric process. Broadcast curvature descends through the tier ladder, accumulating inherited structure, amplifying through tensor channels, and ultimately encountering the collapse–inversion boundary. The visible geometry in Euclidean space is the result of this descent. Positive curvature projects smoothly, producing the ordinary sphere; inverted curvature projects only partially, producing the bounded Amsler patch. Together, these projections reveal the dual structure of the hyperspherical pair; one contracting while one expanding whose primitive superposition defines the curvature field observed in the Euclidean frame.

Projection contains multiple anchors through the metric, broadcast, and inherited curvature with each anchor transforming differently across tiers. Because projection modifies these anchors at different rates, the tensor framework cannot rely on a single blanket constant. Tensors must instead be constructed to reflect tier specific interactions between curvature modes, ensuring that the projected geometry in $\mathbb{R}^3$ accurately represents the dimensional structure from which it descends.

\begin{subsection}{Broadcast, Collapse, and Dimensional Descent}

Curvature originates at higher tiers as broadcast curvature; the expansive propagation of curvature outward from a stability axis. As it descends through the dimensional ladder, it undergoes collapse, a geometric contraction that reduces magnitude and stabilizes orientation. Dimensional descent is therefore a repeated broadcast–collapse sequence, tier by tier, until curvature reaches the seed $T_0.$ This cycle encodes the directional structure of curvature inheritance and determines how curvature becomes visible within the Euclidean frame.
\end{subsection}

\begin{subsection}{Inversion and the Formation of Reciprocal Curvature}

When broadcast curvature exceeds the stabilizing capacity of collapse invariants, the descent does not terminate at $T_0.$ Instead, the stability axis enters the inversion regime. Inversion is the attempt of the curvature field to reverse orientation under excess broadcast, producing reciprocal curvature at $T_{-1},$ the dimensional origin of negative curvature. Reciprocal curvature carries the signatures of inversion: exponential geodesic divergence, expanding metric factors, and the inability to embed globally in Euclidean space. These signatures appear classically in the Amsler surface and dimensionally in the inverted hypersphere, revealing their shared geometric origin.
\end{subsection}

\begin{subsection}{Projection of the First Inversion Tier's Curvature into Euclidean Space}

Curvature at $T_{-1}$ cannot be realized globally when projected into $\mathbb{R}^3.$ Projection is a dimensional reduction: the curvature field is forced into a space that cannot fully contain it. The projection is necessarily bounded, forced to fold, and ultimately self intersect. These behaviors are the collapse–inversion boundary expressed geometrically. Classical geometry records them as Hilbert’s obstruction; Dimensional Calculus identifies them as the projection limit of reciprocal curvature. The visible geometry is therefore not an approximation but the exact Euclidean footprint of inversion curvature attempting to stabilize within a positive curvature frame.
\end{subsection}

\begin{subsection}{Hyperspherical Projection and the Amsler Correspondence}

The inverted hypersphere carries the curvature remainder produced by inversion. When projected into Euclidean space, this remainder becomes the Amsler surface. The correspondence is exact: every intrinsic feature of the Amsler patch matches a dimensional feature of the inverted hypersphere. The boundedness of the Amsler surface reflects the projection limit; its folding reflects inversion; its self intersection reflects the curvature remainder attempting to express itself in a space that cannot fully contain it. Projection Mechanics therefore establishes the geometric equivalence between the inverted hypersphere and the Amsler surface, revealing the Amsler patch as the visible portion of the inverted hypersphere’s tiered curvature remainder.
\end{subsection}

\begin{subsection}{Dual Projection: Positive and Inverted Hyperspheres}

Projection mechanics applies equally to the ordinary hypersphere. Positive curvature projects smoothly into Euclidean space, producing the familiar sphere. Inverted curvature projects only partially, producing the bounded Amsler patch. The two projections, the sphere and the Amsler surface, are the Euclidean images of the two hyperspherical structures. One contracts; one expands. Their superposition produces the full curvature field observed in Euclidean geometry.
\end{subsection}

\begin{subsection}{Physical Interpretation: Gravitational Load and Substrate Expansion}

The curvature constants inherited by the inverted hypersphere carry the inverted gravitational load through the increasing exponential degree of the radius, complete with its inverted effects. Positive curvature produces gravitational contraction; its dimensional dual produces substrate expansion. Projection Mechanics therefore reveals the ``yin–yang'' structure of curvature: the ordinary hypersphere contracts space, while the inverted hypersphere pushes outward. Their bounded projections into Euclidean space are the sphere and the Amsler patch. Their interaction defines the curvature dynamics of the physical world.

In the inverted tier, broadcast curvature is absent. The radius is collapsed and cannot propagate curvature, but the broadcast bar is not inert; anchored in two dimensions, it continues to oscillate, rotate, and sweep laterally under the negative $r$ term. As a result, the orientation coefficient assumes a dual role. It regulates the contraction of positive curvature in our tier, and once inverted, it drives the expansion of subspace. The negative $r$ term remains the metric of the inverted hypersphere: it dictates the oscillatory, two dimensional swing of the broadcast bar even in its collapsed state. Because the inverted tier is gamma generated, its coefficients continue to propagate orientation. When these coefficients slow expansion in our tier, they inflate subspace. This dual behavior is the dimensional signature of reverse gamma: orientation becomes the mechanism of expansion when broadcast is no longer available. As the hypersphere flattens, gravity becomes the sole remaining broadcast channel; with the full unsaturated volume of subspace available, its effect compounds by an additional order of magnitude.
\end{subsection}
\end{section}

\begin{section}{Implications for Curvature Dynamics}

The correspondence between the ordinary hypersphere, the inverted hypersphere, and their Euclidean projections has far reaching implications for Curvature Dynamics. The superposed dual hypersphere structure reveals how curvature propagates, collapses, inverts, and stabilizes across tiers, and how these processes produce the geometric and physical phenomena observed in $\mathbb{R}^3.$ This section summarizes the consequences of Projection Mechanics for curvature deficit, overshoot, sub space geometry, global curvature conservation, and higher tier collapse.

\begin{subsection}{Curvature Deficit and Overshoot}

Curvature deficit arises when broadcast curvature collapses toward the seed $T_0$ and loses magnitude. Classical geometry interprets deficit as “flattening” or “loss of curvature”; DC interprets it as the natural contraction of curvature under dimensional descent. Overshoot occurs when broadcast curvature exceeds collapse invariants and forces the stability axis into inversion. The inverted hypersphere is the geometric manifestation of overshoot. Negative curvature is therefore not an anomaly, but the expected remainder produced when curvature overshoots the stabilizing capacity of collapse.
\end{subsection}

\begin{subsection}{Subspace Geometry}

Subspace is the geometric region formed by the projection of inverted curvature into Euclidean space. It is the bounded, folded, self intersecting region represented by the Amsler surface. Subspace is not a separate physical domain; it is the Euclidean image of $T_{-1}$ curvature. Its geometry reflects the collapse–inversion boundary: boundedness from projection limits, folding from inversion, and self intersection from curvature remainder. Subspace is the dimensional dual of positive curvature space, and its interaction with mass produces annihilation. This duality is the geometric analogue of the gravitational ``yin–yang'': positive curvature contracts, inverted curvature expands.
\end{subsection}

\begin{subsection}{Global Curvature Conservation}

The superposed dual hypersphere structure reveals that curvature is globally conserved across tiers. Positive curvature at higher tiers collapses toward $T_0$; inverted curvature at $T_{-1}$ represents the remainder produced when collapse fails. The sphere and the Amsler patch are therefore complementary projections of a single conserved curvature field. Their superposition produces the full curvature distribution observed in Euclidean geometry. Classical curvature conservation laws are special cases of this broader dimensional conservation principle.
\end{subsection}

\begin{subsection}{Extensions to Higher-Dimensional Collapse}

The collapse–inversion boundary identified in three dimensions extends naturally to higher tiers. Beginning at $T_4,$ curvature amplification generates remainder structures that admit no global Euclidean projection. These higher tier remainders organize into layered subspaces exhibiting increasingly intricate nested torsional behavior. The inverted hypersphere constitutes the first such remainder, with subsequent tiers producing progressively richer curvature signatures. The inverted hypersphere is the first such remainder; higher-tier structures produce additional negative-curvature signatures. Projection Mechanics therefore extends naturally to higher-dimensional collapse, revealing a hierarchy of curvature remainders and their Euclidean manifestations.
\end{subsection}
\end{section}

\begin{section}{Implications For Computational Modeling}

Dimensional Calculus doesn’t just reinterpret curvature — it rewrites the computational landscape. Once the dual hypersphere structure is acknowledged, the PDE machinery that normally governs curvature evolution collapses into a drastically simpler form. The negative tier residues, the inversion boundary at $2\pi$, and the broadcast–inversion superposition together eliminate the pathological behaviors that make classical PDE solvers slow, unstable, or memory hungry. What emerges is a curvature engine that is not only easier to compute, but fundamentally more predictive than the Tier 0 differential framework it replaces.

\begin{subsection}{Python PDE Script Simplification}

Classical PDE solvers struggle because they attempt to evolve curvature using only $T_0$ tools — partial derivatives, smoothness assumptions, and Euclidean constraints. Dimensional Calculus replaces this with a tiered curvature ladder whose behavior is algebraic rather than differential. Python scripts that once required iterative solvers, stability checks, and mesh refinement reduce to direct evaluation of curvature polynomials and their inverted residues. The “PDE” becomes a lookup table on the dimensional ladder. Appendix E provides empirical support for this simplification: a Python generated curvature deficit graph, displayed alongside the PDE expression itself, demonstrates the precise region where Euclid’s model cannot measure curvature. The PDE sits inside the graph’s own frame for direct comparison to the Standard Model, revealing the computational bottleneck that DC eliminates.
\end{subsection}

\begin{subsection}{Latency Bottleneck Evaporates}

The latency in numerical PDE modeling comes from repeated attempts to approximate curvature that is trying to leave the manifold. DC eliminates this by giving each tier a closed form curvature expression. No blow ups, no divergence, no runaway gradients — just direct evaluation of the correct tier. The computational bottleneck evaporates because the curvature engine no longer fights the geometry.
\end{subsection}

\begin{subsection}{Memory Capable Curvature}

The dual hypersphere structure introduces a form of geometric “memory”: each tier inherits constants from the one below it, and the inverted hypersphere retains the broadcast residues in reciprocal form. This gives curvature a built in history that PDEs cannot represent. Computationally, this means curvature evolution is not a blind differential march — it is a structured, tier aware progression with memory encoded directly into the constants $C_1 ,C_2 ,C_3 ,$ and $C_4.$ Python can evaluate this instantly.
\end{subsection}

\begin{subsection}{Cubit Is To Curvature As Curvature Is To Qubit}

The ancient cubit is traditionally treated as a primitive unit of length, but under Dimensional Calculus it reveals itself as a collapsed curvature measure. A cubit is not a “distance” in the Euclidean sense; it is the Tier 0 projection of a higher tier curvature frame. When hyperspherical curvature collapses into the seed state $T_0=1,$ the inherited curvature from $T_1$ is flattened into a linear measurement, producing the familiar forearm length unit. In DC terms, a cubit is the geometric shadow of curvature broadcast: a one tier descent of the radius channel, stripped of its inherited curvature and expressed as a linear segment.

This reinterpretation resolves a long standing historical puzzle: why ancient measurement systems behave as though they encode curvature rather than pure length. The cubit corresponds to the $T_1$ broadcast anchor collapsed into Euclid’s seed frame, carrying the curvature deficit predicted by the $\pi$-bic Lemma. In computational modeling, this means that any “length” measured in cubits is implicitly curvature aware — it inherits the memory of the dimensional tier above it. Thus, the cubit is not merely a unit of human scale measurement; it is the first curvature bearing segment produced by dimensional collapse, a linearized curvature quantity that retains the broadcast signature of the hypersphere.

In the DC framework, curvature itself is not a static geometric quantity; it is a tiered state that inherits information from the dimensional ladder. Each tier carries a distinct curvature signature, and the dual hypersphere structure produces a natural superposition: broadcast curvature on one side of the DC baseline and inverted curvature on the other. This pairing forms a geometric analog of a qubit. A qubit is traditionally defined as a two state quantum object capable of superposition. In DC, curvature occupies two hyperspherical states simultaneously — one expanding, one contracting — with Euclid’s plane acting as the collapse boundary that selects the observed projection.

Unlike quantum qubits, which require probabilistic interpretation, DC curvature states are deterministic: the tiered ladder defines exactly how curvature is inherited, inverted, or collapsed. The “superposition” arises from the dual hypersphere geometry itself, not from uncertainty. When curvature collapses into Euclid’s frame, the projection selects one of the two hyperspherical states, producing the familiar $T_0$ linear geometry. In computational modeling, this means curvature can encode binary like state information without stochastic behavior, offering a geometric pathway to qubit like computation grounded entirely in hyperspherical curvature.
\end{subsection}
\end{section}

\begin{section}{Conclusion}

Dimensional Calculus reveals that negative curvature is not an isolated geometric phenomenon, but the dimensional remainder produced when broadcast curvature overshoots collapse invariants and forces the stability axis into inversion. Classical geometry records the intrinsic behavior of this remainder (exponential geodesic divergence, expanding metric factors, and the impossibility of global embedding) without identifying its dimensional origin. Curvature Dynamics provides that origin by tracing curvature propagation, collapse, inversion, and projection across the tiered curvature ladder.

\begin{subsection}{Summary of Findings}

The central result of this work is the identification of a superposed dual hypersphere structure. The ordinary hypersphere and its inverted counterpart arise from a single seed layer and evolve together as a paired curvature system: one contracting; one expanding. Their bounded projections into Euclidean space, the sphere and the Amsler patch, are complementary images of this superposition. All curvature phenomena observed in $\mathbb{R}^3$ follow from this dual structure. From this central finding, several consequences emerge:

\begin{itemize}
    \item Negative curvature is the reciprocal remainder produced by inversion, not an anomaly. Its boundedness, folding, and self intersection reflect the projection limits of $T_{-1}$ curvature.
    \item Curvature deficit arises from collapse toward the seed tier $T_0,$ while curvature overshoot occurs when broadcast curvature exceeds collapse invariants and forces inversion.
    \item Subspace geometry is the Euclidean image of inverted curvature. The Amsler surface is the exact projection of the inverted hypersphere, with every intrinsic feature corresponding to a dimensional feature of the negative tier curvature remainder.
    \item Global curvature conservation emerges naturally: the sphere and the Amsler patch are two bounded projections of a single conserved curvature field generated by the dual hypersphere pair.
\end{itemize}
\end{subsection}

\begin{subsection}{Reinterpretation of Negative Curvature}

Negative curvature is therefore reinterpreted as the dimensional counterpart to positive curvature. The collapse–inversion boundary explains Hilbert’s obstruction, the behavior of hyperbolic embeddings, and the annihilation boundary between mass and substrate. The asymmetric slopes of the Amsler saddle reflect two distinct curvature residues, $r^{-3}$ and $r^{-4},$ corresponding to $T_{-3}$ and $T_{-4}$ of the dimensional ladder. These residues confirm that the Amsler surface is the Euclidean projection of negative tier curvature broadcast constrained to a two dimensional frame.

Taken together, these results establish that superposition is the foundational structure of Curvature Dynamics. A single seed layer gives rise to both curvature branches (one contracting, one expanding) whose interaction defines the geometry of the physical world. The superposed dual hypersphere system is not a derived feature but the primary architecture of dimensional curvature, and its superposition is the mechanism through which all curvature phenomena in $\mathbb{R}^3$ arise.

Classical geometry treats the Amsler surface as a maximally extended hyperbolic patch that fails to embed smoothly in $\mathbb{R}^3.$ Dimensional Calculus reveals a deeper mechanism: the equations that generate the Amsler surface must undergo eight successive dimensional integrations before their curvature structure aligns with the $T_4$ Euclidean domain. These integrations are not optional; they are dictated by the curvature ladder itself. Crossing the inversion threshold at $2\pi$ is a structural event, marking the transition from negative curvature broadcast to positive curvature reception. The dimensional ladder does not describe how negative curvature embeds in Euclid; it describes how negative curvature integrates into $T_4$ curvature, producing the Euclidean compatible remainder that classical geometry observes. Amsler’s surface is therefore not an embedding problem but an integration pathway through the broadcast–collapse engine of the dual hyperspherical structure of reality.
Einstein’s field equations encode only the broadcast side orientation invariant. The familiar coefficient 8π reflects the curvature contribution of the broadcast hypersphere but omits the inverted hypersphere term required by the full dimensional ladder. As a result, the Einstein tensor captures only half of one of the three geometric invariants, leaving the orientation structure incomplete. Dimensional Calculus restores the missing half, yielding the full orientation coefficient on the order of $16\pi.$
\end{subsection}
\end{section}

\begin{section}{The DC Universe Statement}

Dimensional Calculus reveals that the geometric structure of reality is not a single manifold but a superposition of two hyperspherical curvature domains. The broadcast hypersphere generates ascent curvature through $S$, while its inverted counterpart generates descent curvature through $AA.$ Euclid’s plane lies precisely at the DC baseline inversion point; the mirror boundary where curvature transitions between these domains and undergoes a mandatory $90^\circ$ axis inversion. All observable geometric behavior (including Euclidean linearity, hyperbolic asymmetry, curvature concentration, folding phenomena, and the Gibbs overshoot constant) arises from the interaction of these two hyperspheres across the baseline. The universe is therefore the dynamic superposition of a broadcast hypersphere and its inverted dual, with Euclidean geometry emerging as the projection interface between them.
\end{section}

\begin{section}{Future Directions}

The superposed dual hypersphere model suggests a hierarchy of curvature remainders at higher tiers. Tier $T_{-1}$ produces the inverted hypersphere, while $T_3$ reveals its Euclidean projection limits through boundedness, folding, and self intersection. Higher tiers produce additional negative curvature signatures with increasingly complex projection behavior, forming nested subspace layers that extend the collapse–inversion boundary beyond three dimensions. Future work will explore these higher tier structures, their Euclidean manifestations, and their implications for curvature dynamics, global conservation, and physical geometry. The framework presented here provides the empirical anchor for these investigations, establishing the dimensional origin of negative curvature, the $T_3$ boundary where projection fails, and the superposed dual hypersphere structure that governs all curvature behavior in $\mathbb{R}^3.$

In theory, the dual hypersphere structure also hints at a self correcting computational mechanism: a natural latency load distribution that collapses into coherence. Because curvature exists in superposed broadcast and inverted states, the system behaves like a geometric load balancer — one tier expanding, one tier contracting, each compensating for the other. This mirrors the classical processor “pin” model for data retrieval, but replaces electronic state checking with curvature state selection. In DC terms, coherence is not an emergent property of computation; it is the deterministic outcome of hyperspherical superposition.
\end{section}

\begin{section}{Ethics Statement}

This work develops a geometric framework for curvature behavior across hyperspherical and inverted domains. While the implications of Dimensional Calculus extend to fields such as physics, geometry, and computational modeling, all results presented here are theoretical. No part of this framework is intended for harmful applications, and all interpretations are offered solely for constructive scientific inquiry.
\end{section}

\begin{section}{AI-Assistance Disclosure}

Generative AI (Microsoft Copilot) was used for text refinement, structural editing, and accessibility support. All mathematical results, derivations, figures, and computational interpretations were developed solely by the author.
\end{section}

\begin{section}{Python Computational Methods}

Python was used to generate tier curve plots, inversion boundary behavior, and subspace asymptote visualizations. All scripts were written by the author and executed locally. These computational results serve as empirical illustrations of the curvature ladder and its projection limits.
\end{section}

\begin{section}{Funding}

This research received no external funding.
\end{section}

\begin{section}{Conflicts of Interest}

The author declares no conflicts of interest.
\end{section}

\begin{section}{Data Availability}

All data used in this work are theoretical constructions developed by the author. Python scripts and supplementary figures are available upon request.
\end{section}

\begin{thebibliography}{99}

\bibitem{Amsler1955}
M.~Amsler.
\newblock Sur les surfaces à courbure négative.
\newblock \emph{Bulletin de la Société Mathématique de France}, 83:1--48, 1955.
\newblock DOI: \href{https://doi.org/10.24033/bsmf.1461}{10.24033/bsmf.1461}.

\bibitem{Schoenman2026a}
A.~Schoenman.
\newblock \emph{Dimensional Calculus: A Guide to Mapping Reality Through Curvature}.
\newblock Columbus, Ohio, 2026.
\newblock DOI: \href{https://doi.org/10.5281/zenodo.21616655}{10.5281/zenodo.21616655}.

\bibitem{Schoenman2026b}
A.~Schoenman.
\newblock \emph{Axioms of Curvature: Revealing and Correcting Euclid's Hidden Assumption}.
\newblock Version 3.0, 2026.
\newblock DOI: \href{https://doi.org/10.5281/zenodo.22073321}{10.5281/zenodo.22073321}.

\bibitem{Schoenman2026c}
A.~Schoenman.
\newblock \emph{The Gibbs Overshoot Mapped Through Curvature Dynamics: The Simplification Power of the Hypersphere}.
\newblock Columbus, Ohio, 2026.
\newblock DOI: \href{https://doi.org/10.5281/zenodo.21982855}{10.5281/zenodo.21982855}.

\bibitem{Schoenman2026d}
A.~Schoenman.
\newblock \emph{Axioms of Broadcast: The Geometric Emergence of Gravity, Collapse, and Cosmogenesis}.
\newblock Columbus, Ohio, 2026.
\newblock DOI: \href{https://doi.org/10.5281/zenodo.22085513}{10.5281/zenodo.22085513}.

\end{thebibliography}
\end{document}